\documentclass[
 aps,
 prd,
 reprint,
 amsmath,
 amssymb,
 nofootinbib,
 superscriptaddress
]{revtex4-2}

\usepackage{mathtools}
\usepackage{amsthm}
\usepackage{bm}
\usepackage{mathrsfs}
\usepackage{booktabs}
\usepackage{enumitem}
\usepackage{hyperref}
\usepackage[nameinlink,capitalize]{cleveref}

\hypersetup{
 colorlinks=true,
 linkcolor=blue,
 citecolor=blue,
 urlcolor=blue
}

\newtheorem{theorem}{Theorem}[section]
\newtheorem{proposition}[theorem]{Proposition}
\newtheorem{lemma}[theorem]{Lemma}
\newtheorem{corollary}[theorem]{Corollary}
\theoremstyle{definition}
\newtheorem{definition}[theorem]{Definition}
\theoremstyle{remark}
\newtheorem{remark}[theorem]{Remark}

\newcommand{\Sthree}{S^3_R}
\newcommand{\PP}{\mathbb P}
\newcommand{\cL}{\mathcal L}
\newcommand{\cG}{\mathcal G}
\newcommand{\cM}{\mathscr M}
\newcommand{\cS}{\mathscr S}
\newcommand{\cZ}{\mathcal Z}
\newcommand{\cE}{\mathcal E}
\newcommand{\cK}{\mathfrak K}
\newcommand{\cA}{\mathscr A}
\newcommand{\cR}{\mathscr R}
\newcommand{\cB}{\mathscr B}
\newcommand{\cC}{\mathscr C}
\newcommand{\cN}{\mathscr N}
\newcommand{\curl}{\operatorname{curl}}
\newcommand{\divg}{\operatorname{div}_g}
\newcommand{\gradg}{\operatorname{grad}_g}
\newcommand{\Def}{\operatorname{Def}}
\newcommand{\Id}{\mathrm{Id}}
\newcommand{\Gr}{\operatorname{Gr}}
\newcommand{\Vol}{\operatorname{Vol}}
\newcommand{\dd}{\,\mathrm d}
\newcommand{\HS}{\mathrm{HS}}
\newcommand{\tr}{\operatorname{tr}}
\newcommand{\op}{\mathrm{op}}
\newcommand{\ad}{\operatorname{ad}}

\begin{document}

\title{Geometric Reduction, Killing-Response Reconstruction, and Spectral Tomography for Navier--Stokes Flow on the Round Three-Sphere}

\author{Boris~Batenin}
\email{batenin@web.de}
\affiliation{Independent Researcher, Berlin, Germany}

\author{Andrei~Preece}
\email{andre.preece@web.de}
\affiliation{Independent Researcher, Berlin, Germany}

\begin{abstract}
We develop a unified geometric and spectral theory for incompressible
Navier--Stokes flow on the round three-sphere.  The common hinge is the
Killing--Beltrami interaction on the signed curl decomposition.  First, the
projected quadratic convection vanishes on every single signed curl shell,
so the full Grassmannian of its linear subspaces consists of exact
finite-dimensional Navier--Stokes reductions.  We study the resulting
moduli under the $SO(4)$ isometry action, derive a real-algebraic leakage
functional for multi-shell reductions, and obtain its infinitesimal
deformation equation.  Killing drift transports these reductions by a
projector Lax equation; for constant drift we give an exact Grassmannian
geodesic criterion, diagonalize the shell generator, derive the finite
Diophantine resonance condition, and prove determinant-one shell transport.
Second, the six convective responses to Killing generators factor the
Stokes operator.  After $A^{-1/2}$ normalization they form a Parseval
operator-valued frame.  Either chiral triplet is uniformly stable, two
orthonormal same-chirality probes determine the entire non-Killing velocity,
and a single probe is invertible exactly on odd signed-curl levels.  Gram
compression yields a Casimir response tensor with exact chiral trace
splitting.  Stokes weighting produces a positive Hankel moment hierarchy,
while curl weighting produces a positive matrix-valued spectral measure.
Because the signed curl spectrum lies on an integer lattice, its
operator-valued characteristic function is periodic; Fourier inversion
recovers every signed-shell response tensor, and its scalar trace recovers
all signed-shell energies for arbitrary finite-dissipation states.  Under
linear Stokes evolution this characteristic function satisfies an exact
matrix heat equation on the response circle.  Combining the reduction and
response branches further yields Clebsch--Gordan parity and locality rules
for nonlinear shell interactions, one- and two-probe certificates for the
full leakage variety, a master response-circle functional calculus, and a
Fisher--Rao spectral-selection law whose nonlinear extension separates
viscous gradient drift from conservative shell transfer.
\end{abstract}
\maketitle

\section{Introduction}

The round three-sphere is a particularly rigid background for incompressible
hydrodynamics: its isometry group is \(SO(4)\), its Killing algebra splits
into two commuting \(\mathfrak{su}(2)\) factors, and the divergence-free
spectrum admits a signed curl decomposition.  With deformation-tensor
viscosity these structures interact directly with the projected
Navier--Stokes nonlinearity.

The purpose of this article is to treat two consequences of the same
Killing--Beltrami algebra within one theorem chain.  The first is a
\emph{reduction problem}: which finite-dimensional subspaces are exact
Navier--Stokes systems, how are they organized geometrically, and how do
Killing drifts move them?  The second is an \emph{information problem}: how
much of a velocity field can be recovered from its convective responses to
the six Killing generators?

The reduction problem begins with a complete cancellation on a single
signed curl shell.  Since the Stokes operator is scalar there and the
polarized projected convection vanishes, every linear subspace of that shell
is an exact pressure-linear reduction.  Thus exact reductions occur in
Grassmannian families rather than as isolated modes.  The physical
\(SO(4)\) action produces quotient moduli, while cross-shell interactions
lead to a real-algebraic zero-leakage locus on products of Grassmannians.
Killing drift gives a natural nonabelian transport on these moduli.  Its
projector satisfies a Lax equation, but a constant drift orbit is geodesic
only under an additional block-compatibility condition.  The same
representation theory gives the exact fixed-drift spectrum and its
resonance arithmetic.

The information problem begins from the quadratic Casimir factorization of
the Stokes operator by six Killing-response operators.  After Stokes
normalization these responses form a Parseval operator-valued frame.  The
redundancy is explicit: either chiral triplet remains a stable frame, any
two orthonormal probes in one chiral factor are globally injective, and the
one-probe problem has an exact odd/even shell parity law.

Compressing the response fields to their \(6\times6\) Gram operator yields
a Casimir response geometry.  Applying powers of the Stokes operator
produces a positive moment hierarchy whose traces are the Stokes spectral
moments.  Applying powers of curl instead retains chirality and yields a
positive matrix-valued measure on the signed curl spectrum.  Since
\(\mu_{k,\sigma}=\sigma(k+2)/R\), the associated unitary curl group is
periodic with period \(2\pi R\); consequently a single operator-valued
function on this response circle admits exact Fourier inversion.

The individual ingredients have established relatives: invariant helical
subspaces and special exact solutions occur in Navier--Stokes theory
\cite{MahalovTitiLeibovich1990,BogoyavlenskijFuchssteiner2005}; Killing
fields in deformation-tensor Navier--Stokes equations were studied by
Samavaki and Tuomela \cite{SamavakiTuomela2020}; Grassmann geometry,
operator-valued frames, Lie-group Riesz transforms, and matrix-valued moment
problems have independent abstract theories
\cite{Bendokat2024,Sun2006,KaftalLarsonZhang2009,
BoggarapuThangavelu2016,KimseyTrachana2022}.  The contribution here is the
exact realization of these structures by one \(SO(4)\)-equivariant
Navier--Stokes system and the explicit theorems linking reduction geometry,
minimal observability, and spectral tomography.  No historical priority is
inferred from the literature search; the claims are the proved statements
below.



\subsection*{Relation to the three public Zenodo preprints}

The present manuscript belongs to a three-stage public preprint program.

The earliest work, Ref.~\cite{BateninPreece2026}, was publicly deposited on Zenodo on
September 20, 2026.  It introduced the general finite-dimensional closure
criterion and the zero-defect viewpoint, developed the maximal-torus,
Killing, and signed-Beltrami sectors, established spectral-escape/projector
bounds, and kept the forced singular-transfer construction explicitly
conditional.

The second work, Ref.~\cite{BateninPreeceKB2026}, was publicly deposited on Zenodo on
September 21, 2026.  It isolated the exact Killing--Beltrami intertwining
law, shell-diagonal Killing precession, the associated connection and
curvature, Killing-saturated rigidity, the quadratic Casimir factorization,
a finite toric classification, and sharp homogeneous-projector norms.  That
paper is not superseded by the unified manuscript: several of these
classification and homogeneous-bundle results are not reproduced below.

The third work, Ref.~\cite{BateninPreeceUnified2026}, was also publicly deposited on
September 21, 2026 and is the immediate Version 1.0 baseline for the unified
reduction/response theory.  It developed Grassmannian reduction moduli,
multi-shell leakage geometry, Killing-precession geometry, lossless
Killing-response reconstruction, Casimir response moments, and
curl-Fourier tomography.

The additional results collected in Sec.~\ref{N:sec:synthesis} were derived
subsequently.  In particular, the Clebsch--Gordan interaction selection rule
there answers a limitation explicitly left open in
Ref.~\cite{BateninPreeceKB2026}: that paper did not infer a chiral shell
selection rule without an explicit Clebsch--Gordan analysis.  The new
section should therefore be read as a subsequent extension of the three
public preprints, not as a retrospective claim about their original
contents.

\section{Geometric and spectral setting}

Let $(\Sthree,g)$ be the round three-sphere of radius $R$, with
\[
        V_R=\Vol(\Sthree)=2\pi^2R^3.
\]
We use the deformation-tensor viscosity
\[
\begin{aligned}
 \partial_tu+\nabla_u u-\nu\cL_g u+\gradg p&=f,\\
 \divg u&=0,\\
 \cL_g u&=2\,\operatorname{div}_g\Def_g u.
\end{aligned}
\]
Let $\PP$ denote the Leray projection and put
\[
\begin{aligned}
 A&=-\PP\cL_g,\\
 B(u,v)&=\PP\nabla_u v,\\
 Q(u,v)&=\frac12\{B(u,v)+B(v,u)\}.
\end{aligned}
\]
Then
\begin{equation}
       \partial_tu+\nu Au+Q(u,u)=\PP f.
       \label{G:eq:NS}
\end{equation}

The divergence-free $L^2$ space decomposes as
\begin{equation}
 L^2_\sigma(T\Sthree)
 =
 \widehat{\bigoplus}_{k\ge0,\ \sigma=\pm} E_{k,\sigma},
 \label{G:eq:signeddecomp}
\end{equation}
where
\begin{align}
 \curl u &= \mu_{k,\sigma}u,
 &
 \mu_{k,\sigma}&=\sigma\frac{k+2}{R},
 \label{G:eq:curlspec}\\
 Au &= \kappa_k u,
 &
 \kappa_k&=\frac{k(k+4)}{R^2},
 \label{G:eq:stokespec}\\
 \dim E_{k,\sigma}
 &=d_k=(k+1)(k+3).
 \label{G:eq:dim}
\end{align}
These formulas are consistent with the standard spectrum of differential
forms on spheres \cite{IkedaTaniguchi1978} and with the signed-shell
conventions of Ref.~\cite{BateninPreece2026}.

We first record the cancellation that drives everything that follows.

\begin{lemma}[Polarized single-shell cancellation]
\label{G:lem:shellcancel}
For all $u,v\in E_{k,\sigma}$,
\begin{equation}
          Q(u,v)=0.
          \label{G:eq:shellcancel}
\end{equation}
\end{lemma}

\begin{proof}
On an oriented Riemannian three-manifold,
\[
 \nabla_uv+\nabla_vu
 =
 \gradg\langle u,v\rangle
 -u\times\curl v-v\times\curl u.
\]
If $u$ and $v$ belong to the same signed curl shell, then
$\curl u=\mu u$ and $\curl v=\mu v$ with the same scalar $\mu$.  Hence
\[
 -u\times\curl v-v\times\curl u
 =
 -\mu u\times v-\mu v\times u=0.
\]
Thus $\nabla_uv+\nabla_vu$ is a gradient, and its Leray projection vanishes.
\end{proof}

\begin{remark}
For $u=v$, \cref{G:lem:shellcancel} is the usual Beltrami observation that
$\nabla_u u$ is a gradient.  The polarized statement is stronger: it removes
the interaction between \emph{any two} vectors in the same signed shell.
\end{remark}


\section{Exact-reduction moduli and leakage geometry}

\subsection{The Grassmannian exact-reduction theorem}

Fix $(k,\sigma)$ and abbreviate
\[
          E=E_{k,\sigma},
          \qquad
          D=\dim E=d_k.
\]
All Grassmannians below are real.

\begin{theorem}[Grassmannian family of exact reductions]
\label{G:thm:grassmann}
Let $1\le d\le D$.  Every point
\[
          L\in\Gr(d,E)
\]
is an exact $d$-dimensional reducing space for \eqref{G:eq:NS}:
\begin{equation}
          AL\subset L,
          \qquad
          Q(L,L)=0.
          \label{G:eq:exactL}
\end{equation}
For unforced dynamics and $u_0\in L$,
\begin{equation}
          u(t)=e^{-\nu\kappa_k t}u_0.
          \label{G:eq:fiberflow}
\end{equation}
If $\PP f(t)\in L$ for all $t$, then
\begin{equation}
 u(t)
 =
 e^{-\nu\kappa_k t}u_0
 +
 \int_0^t
 e^{-\nu\kappa_k(t-s)}\PP f(s)\,\dd s.
 \label{G:eq:forcedfiber}
\end{equation}
\end{theorem}

\begin{proof}
Equation \eqref{G:eq:stokespec} gives
$A|_E=\kappa_k\Id_E$, hence $AL\subset L$ for every subspace $L\subset E$.
Equation \eqref{G:eq:shellcancel} gives $Q(L,L)=0$.  Restricting
\eqref{G:eq:NS} to $L$ therefore leaves the linear scalar equation
\[
        \partial_tu+\nu\kappa_k u=\PP f.
\]
The formulas follow by variation of constants.
\end{proof}

\begin{corollary}[No finite low-dimensional atlas]
\label{G:cor:nofiniteatlas}
For every $d\ge1$ there are infinitely many exact $d$-dimensional linear
Navier--Stokes reductions on $\Sthree$.  For every $1\le d\le8$ there are
continuous families already inside a single signed curl shell.
\end{corollary}

\begin{proof}
Choose $k$ so that $d_k>d$.  Then
$\Gr(d,E_{k,\sigma})$ is a positive-dimensional compact manifold of
dimension $d(d_k-d)$.
\end{proof}

\begin{remark}
The theorem concerns exact \emph{linear} invariant subspaces.  It is
different from the local nonlinear spectral submanifolds used in modern
model reduction \cite{Buza2023}.  Here invariance is global on the subspace
and no perturbative graph construction is required.
\end{remark}

\subsection{Tautological Navier--Stokes dynamics}

Let
\[
 \cG_{D,d}:=\Gr(d,E)
\]
and let
\[
 \pi:\cS_{D,d}\longrightarrow\cG_{D,d}
\]
be the tautological bundle,
\[
 \cS_{D,d}
 =
 \{(L,u)\in\cG_{D,d}\times E:\ u\in L\}.
\]

\begin{theorem}[Universal tautological semigroup]
\label{G:thm:tautological}
For each $t\ge0$, define
\begin{equation}
 \Phi_t:\cS_{D,d}\longrightarrow\cS_{D,d},
 \qquad
 \Phi_t(L,u)
 =
 \bigl(L,e^{-\nu\kappa_k t}u\bigr).
 \label{G:eq:tautflow}
\end{equation}
Then $(\Phi_t)_{t\ge0}$ is a smooth bundle semigroup over the identity on
$\cG_{D,d}$, and its restriction to every fiber is exactly the unforced
Navier--Stokes flow on that reduction.

The construction is $SO(4)$-equivariant.
\end{theorem}

\begin{proof}
Fiber preservation and equality with the Navier--Stokes flow are
\cref{G:thm:grassmann}.  The formula is smooth and satisfies
$\Phi_{t+s}=\Phi_t\circ\Phi_s$.  Since the $SO(4)$ action preserves the round
metric, curl shells, and the Stokes operator,
\[
 g\cdot\Phi_t(L,u)
 =
 (gL,e^{-\nu\kappa_k t}gu)
 =
 \Phi_t(gL,gu).
\]
\end{proof}

\begin{remark}
The base point $L$ labels the exact reduced model; the fiber variable is the
fluid state within that model.  Thus the tautological bundle packages all
$d$-dimensional same-shell exact reductions into one geometric dynamical
object.
\end{remark}

\subsection{Isometry moduli and orbit dimension}

The isometry group $G=SO(4)$ acts orthogonally on $E_{k,\sigma}$ and hence
smoothly on $\Gr(d,E_{k,\sigma})$.

\begin{definition}
The moduli space of $d$-dimensional same-shell reductions up to isometry is
\begin{equation}
 \cM_{k,\sigma,d}
 :=
 \Gr(d,E_{k,\sigma})/SO(4).
 \label{G:eq:moduli}
\end{equation}
\end{definition}

\begin{theorem}[Compact quotient and dimension lower bound]
\label{G:thm:moddim}
The quotient $\cM_{k,\sigma,d}$ is compact and Hausdorff.  Its principal
orbit stratum has dimension
\begin{equation}
 d(d_k-d)-r_{\mathrm{pr}},
 \label{G:eq:principalmoddim}
\end{equation}
where $r_{\mathrm{pr}}\le6$ is the dimension of a principal $SO(4)$ orbit.
Consequently,
\begin{equation}
 \boxed{
 \dim\cM_{k,\sigma,d}
 \ge
 d(d_k-d)-6.
 }
 \label{G:eq:modlower}
\end{equation}
In particular, if $d(d_k-d)>6$, there is a positive-dimensional continuum
of pairwise isometrically inequivalent exact reductions.
\end{theorem}

\begin{proof}
The real Grassmannian is compact and $SO(4)$ is compact, so the orbit space
is compact Hausdorff.  By the principal orbit type theorem, a dense open
subset of the Grassmannian consists of orbits of one principal type.  If
their dimension is $r_{\mathrm{pr}}$, then the quotient of this principal
stratum has dimension
\[
      \dim\Gr(d,D)-r_{\mathrm{pr}}
      =
      d(D-d)-r_{\mathrm{pr}}.
\]
Since $\dim SO(4)=6$, one has $r_{\mathrm{pr}}\le6$.
\end{proof}

\begin{corollary}[Positive-dimensional moduli in the first non-Killing shell]
\label{G:cor:firstshellmoduli}
For $k=1$, one has $d_1=8$.  Hence for every proper dimension
$1\le d\le7$,
\[
        \dim\cM_{1,\sigma,d}\ge d(8-d)-6>0.
\]
Thus even the first non-Killing signed curl shell contains continuous
isometry moduli of exact reductions in every proper dimension.
\end{corollary}

\begin{remark}
The quotient need not be a smooth manifold globally because stabilizer type
can jump.  The theorem uses only the principal stratum, so no global
freeness assumption is made.
\end{remark}

\subsection{A concentration landscape on the Grassmannian}

Let $L$ be any $d$-dimensional subspace of smooth vector fields on $\Sthree$,
and let $\Pi_L$ be its $L^2$-orthogonal projector.  Choose an orthonormal
basis $\{\phi_j\}_{j=1}^d$ of $L$.  The projector kernel is
\begin{equation}
 K_L(x,y)
 =
 \sum_{j=1}^d
 \phi_j(x)\otimes\phi_j(y)^\flat.
 \label{G:eq:kernel}
\end{equation}
At the diagonal $K_L(x,x)$ is a positive semidefinite endomorphism of the
three-dimensional tangent space.

\begin{lemma}[Evaluation formula]
\label{G:lem:evaluation}
For the evaluation map
$E_x:L\to T_x\Sthree$, one has
\[
 K_L(x,x)=E_xE_x^*,
\]
and therefore
\begin{equation}
 \|\Pi_L\|_{2\to\infty}^2
 =
 \sup_{x\in\Sthree}
 \lambda_{\max}\bigl(K_L(x,x)\bigr).
 \label{G:eq:L2kernel}
\end{equation}
Moreover,
\begin{equation}
 \int_{\Sthree}\tr K_L(x,x)\,\dd V=d.
 \label{G:eq:traceidentity}
\end{equation}
\end{lemma}

\begin{proof}
The kernel identity follows directly from an orthonormal basis.  The
$L^2\to L^\infty$ norm is the supremum of the evaluation norms, and
$\|E_x\|^2=\lambda_{\max}(E_xE_x^*)$.  Finally,
\[
 \int\tr K_L(x,x)\dd V
 =
 \sum_{j=1}^d\|\phi_j\|_2^2=d.
\]
\end{proof}

\begin{theorem}[Universal concentration floor]
\label{G:thm:concentrationfloor}
For every $d$-dimensional subspace $L$ of smooth vector fields on $\Sthree$,
\begin{align}
 \boxed{
 \|\Pi_L\|_{2\to\infty}
 \ge
 \sqrt{\frac{d}{3V_R}},
 }
 \label{G:eq:floor2}\\
 \boxed{
 \|\Pi_L\|_{1\to\infty}
 \ge
 \frac{d}{3V_R}.
 }
 \label{G:eq:floor1}
\end{align}
In particular these inequalities hold for every point of every exact
Grassmannian family $\Gr(d,E_{k,\sigma})$.
\end{theorem}

\begin{proof}
By \eqref{G:eq:traceidentity}, there exists $x$ with
\[
       \tr K_L(x,x)\ge\frac{d}{V_R}.
\]
Since $K_L(x,x)$ acts on a three-dimensional space,
\[
 \lambda_{\max}\bigl(K_L(x,x)\bigr)
 \ge
 \frac13\tr K_L(x,x)
 \ge
 \frac{d}{3V_R}.
\]
Equation \eqref{G:eq:floor2} follows from \eqref{G:eq:L2kernel}.

For the $L^1\to L^\infty$ norm, the continuous finite-rank kernel satisfies
\[
 \|\Pi_L\|_{1\to\infty}
 =
 \sup_{x,y}\|K_L(x,y)\|_{\op}.
\]
The diagonal is included in the supremum, so the same point $x$ gives
\eqref{G:eq:floor1}.
\end{proof}

\begin{definition}
For $L\in\Gr(d,E_{k,\sigma})$, define the dimensionless concentration
defects
\begin{equation}
 \mathcal C_2(L)
 :=
 \frac{3V_R}{d}\|\Pi_L\|_{2\to\infty}^2,
 \qquad
 \mathcal C_1(L)
 :=
 \frac{3V_R}{d}\|\Pi_L\|_{1\to\infty}.
 \label{G:eq:defects}
\end{equation}
Then $\mathcal C_1,\mathcal C_2\ge1$.
\end{definition}

These functions are invariant under the physical $SO(4)$ action and
therefore descend to the isometry moduli space \eqref{G:eq:moduli}.  They
provide quantitative functions on the moduli of exact reductions.

\subsection{Random Grassmannian reductions}

We now equip $\Gr(d,E)$ with its normalized $O(E)$-invariant probability
measure.  This ensemble is not a physical symmetry average; it is the
canonical uniform measure on the parameter space of exact reductions.

Let $P_L:E\to E$ denote the orthogonal projector onto $L$.

\begin{lemma}[First moment of a random subspace projector]
\label{G:lem:firstmoment}
If $L$ is uniformly distributed in $\Gr(d,E)$ and $D=\dim E$, then
\begin{equation}
       \mathbb E\,P_L=\frac{d}{D}\Id_E.
       \label{G:eq:firstmoment}
\end{equation}
\end{lemma}

\begin{proof}
The expectation commutes with every orthogonal transformation of $E$, so it
is a scalar multiple of the identity.  Taking the trace gives the scalar
$d/D$.
\end{proof}

Let $K_E(x,y)$ be the orthogonal projector kernel of the full shell.  By
$SO(4)$ homogeneity and irreducibility of the tangent isotropy,
\begin{equation}
       K_E(x,x)
       =
       \frac{D}{3V_R}\Id_{T_x\Sthree}.
       \label{G:eq:fullshelldiagonal}
\end{equation}

\begin{theorem}[Ensemble isotropy]
\label{G:thm:ensembleisotropy}
For a uniformly random $L\in\Gr(d,E_{k,\sigma})$,
\begin{align}
 \mathbb E\,K_L(x,y)
 &=
 \frac{d}{D}K_E(x,y),
 \label{G:eq:meankernel}\\
 \boxed{
 \mathbb E\,K_L(x,x)
 =
 \frac{d}{3V_R}\Id_{T_x\Sthree}.
 }
 \label{G:eq:meandiagonal}
\end{align}
Thus the ensemble of exact reductions is pointwise isotropic in expectation,
even though an individual reduction need not be homogeneous.
\end{theorem}

\begin{proof}
Write $\operatorname{ev}_x:E\to T_x\Sthree$ for evaluation.  Then
\[
 K_L(x,y)
 =
 \operatorname{ev}_x P_L\operatorname{ev}_y^*.
\]
Take expectation and use \eqref{G:eq:firstmoment}; the diagonal formula follows
from \eqref{G:eq:fullshelldiagonal}.
\end{proof}

The second moment is also explicit.

\begin{theorem}[Exact pointwise projector variance]
\label{G:thm:variance}
Let $D=d_k$ and $1\le d\le D$.  For every $x\in\Sthree$,
\begin{equation}
 \boxed{
 \mathbb E
 \left\|
 K_L(x,x)
 -
 \frac{d}{3V_R}\Id
 \right\|_{\HS}^2
 =
 \frac{
 2d(D-d)(2D-1)
 }{
 3V_R^2(D-1)(D+2)
 }.
 }
 \label{G:eq:variance}
\end{equation}
The right-hand side is independent of $x$.
\end{theorem}

\begin{proof}
Fix an orthonormal basis of $E$ and write $P=P_L$.  Orthogonal invariance
implies
\begin{equation}
 \mathbb E(P_{ij}P_{kl})
 =
 a\,\delta_{ij}\delta_{kl}
 +
 b\bigl(
 \delta_{ik}\delta_{jl}
 +
 \delta_{il}\delta_{jk}
 \bigr).
 \label{G:eq:secondtensor}
\end{equation}
The identities
$\tr P=d$ and $\tr P^2=d$ give
\begin{align}
 a&=
 \frac{d\{d(D+1)-2\}}
 {D(D-1)(D+2)},
 \label{G:eq:a}\\
 b&=
 \frac{d(D-d)}
 {D(D-1)(D+2)}.
 \label{G:eq:b}
\end{align}

Let $T=\operatorname{ev}_x:E\to T_x\Sthree$ and
$H=T^*T$.  Then
$K_L(x,x)=TPT^*$ and, by \eqref{G:eq:fullshelldiagonal},
\begin{equation}
 \tr H=\frac{D}{V_R},
 \qquad
 \tr H^2=\frac{D^2}{3V_R^2}.
 \label{G:eq:Hmoments}
\end{equation}
Contracting \eqref{G:eq:secondtensor} yields
\begin{equation}
 \mathbb E\,\tr(TPT^*)^2
 =
 (a+b)\tr H^2+b(\tr H)^2.
 \label{G:eq:Ksecond}
\end{equation}
Subtract
\[
 \left\|
 \frac{d}{3V_R}\Id
 \right\|_{\HS}^2
 =
 \frac{d^2}{3V_R^2},
\]
insert \eqref{G:eq:a}--\eqref{G:eq:Hmoments}, and simplify.  This gives
\eqref{G:eq:variance}.
\end{proof}

\begin{remark}
For a fixed proportion $d/D$, the absolute variance grows only linearly in
$D$, whereas the squared size of the mean diagonal kernel grows
quadratically.  Thus high-dimensional random exact reductions become
relatively isotropic at a fixed point.
\end{remark}

\subsection{Multi-shell leakage as a real-algebraic moduli problem}

The single-shell Grassmannian is maximally degenerate: every point is exact.
Cross-shell interaction creates the genuine classification problem.

Let
\[
 \mathcal A=\{\alpha_1,\ldots,\alpha_m\},
 \qquad
 E_{\alpha_j}=E_{k_j,\sigma_j},
\]
be finitely many signed shells and set
\[
        V_{\mathcal A}
        =
        \bigoplus_{\alpha\in\mathcal A}E_\alpha.
\]
Fix a dimension vector
$\mathbf d=(d_\alpha)_{\alpha\in\mathcal A}$ with
$0\le d_\alpha\le\dim E_\alpha$, and define
\begin{equation}
 \cG_{\mathcal A,\mathbf d}
 :=
 \prod_{\alpha\in\mathcal A}
 \Gr(d_\alpha,E_\alpha).
 \label{G:eq:productgrass}
\end{equation}
A point
$\mathbf L=(L_\alpha)_\alpha$
determines
\[
        L(\mathbf L)
        =
        \bigoplus_{\alpha\in\mathcal A}L_\alpha
\]
and its orthogonal projector
\[
        P_{\mathbf L}
        =
        \bigoplus_{\alpha\in\mathcal A}P_{L_\alpha},
\]
extended by zero on $V_{\mathcal A}^\perp$.

Because $A$ is scalar on each shell,
\[
        A L(\mathbf L)\subset L(\mathbf L)
\]
for every point of \eqref{G:eq:productgrass}.  Only quadratic leakage remains.

\begin{definition}[Grassmannian leakage functional]
Choose any orthonormal basis $\{e_i\}_{i=1}^N$ of
$V_{\mathcal A}$.  Define
\begin{equation}
 \Lambda_{\mathcal A,\mathbf d}(\mathbf L)
 :=
 \sum_{i,j=1}^N
 \left\|
 (I-P_{\mathbf L})
 Q(P_{\mathbf L}e_i,P_{\mathbf L}e_j)
 \right\|_2^2.
 \label{G:eq:leakage}
\end{equation}
\end{definition}

This is the squared Hilbert--Schmidt norm of the bilinear map
\[
 (I-P_{\mathbf L})Q
 \big|_{L(\mathbf L)\times L(\mathbf L)},
\]
so it is independent of the chosen orthonormal basis.

\begin{theorem}[Leakage variety]
\label{G:thm:leakagevariety}
The function
$\Lambda_{\mathcal A,\mathbf d}$ has the following properties.

\begin{enumerate}[label=(\roman*)]
\item
It is nonnegative, continuous, real analytic, and $SO(4)$-invariant.

\item
In the orthogonal-projector realization
\[
 \Gr(d,E)
 \simeq
 \{P\in\operatorname{End}(E):
 P^2=P,\ P^*=P,\ \tr P=d\},
\]
it is the restriction of a polynomial function of the projector entries.

\item
Its zero set
\begin{equation}
 \cZ_{\mathcal A,\mathbf d}
 :=
 \{\mathbf L\in\cG_{\mathcal A,\mathbf d}:
 \Lambda_{\mathcal A,\mathbf d}(\mathbf L)=0\}
 \label{G:eq:Z}
\end{equation}
is exactly the set of exact reducing spaces with the prescribed spectral
dimension vector:
\begin{equation}
 \mathbf L\in\cZ_{\mathcal A,\mathbf d}
 \quad\Longleftrightarrow\quad
 Q(L(\mathbf L),L(\mathbf L))
 \subset L(\mathbf L).
 \label{G:eq:Zcriterion}
\end{equation}

\item
$\cZ_{\mathcal A,\mathbf d}$ is a compact $SO(4)$-invariant real-algebraic
subset of the compact product Grassmannian.
\end{enumerate}
\end{theorem}

\begin{proof}
Nonnegativity and continuity are immediate.  The Hilbert--Schmidt expression
is independent of basis.  Equivariance of the Navier--Stokes bilinear map
under isometries gives $SO(4)$ invariance.

For algebraicity, use orthogonal projector coordinates.  The vectors
$P_{\mathbf L}e_i$ depend linearly on the entries of $P_{\mathbf L}$,
the bilinear map $Q$ is fixed, and the final application of
$I-P_{\mathbf L}$ is again linear in the projector.  Since
$Q(V_{\mathcal A},V_{\mathcal A})$ spans a finite-dimensional space, the
squared $L^2$ norm in \eqref{G:eq:leakage} is a finite polynomial expression in
those entries.

The sum of squares vanishes exactly when every leakage vector vanishes,
which is precisely
\[
      Q(L(\mathbf L),L(\mathbf L))
      \subset L(\mathbf L).
\]
Stokes invariance is automatic from the shell decomposition, proving
\eqref{G:eq:Zcriterion}.  Compactness follows because the zero set is closed
inside a compact product of Grassmannians.
\end{proof}

\begin{corollary}[Single-shell component]
\label{G:cor:singlecomponent}
If $\mathcal A=\{\alpha\}$ consists of one signed shell, then
\[
        \Lambda_{\mathcal A,d}\equiv0,
        \qquad
        \cZ_{\mathcal A,d}
        =
        \Gr(d,E_\alpha).
\]
\end{corollary}

\begin{corollary}[Complete-shell leakage is a vertex specialization]
\label{G:cor:vertex}
Suppose each $d_\alpha$ is either $0$ or $\dim E_\alpha$.  Let $S$ be the
set of selected complete shells.  Then
\begin{equation}
 \Lambda_{\mathcal A,\mathbf d}
 =
 \sum_{\alpha,\beta\in S}
 \sum_{\gamma\notin S}
 \left\|
 \Pi_\gamma Q
 \big|_{E_\alpha\times E_\beta}
 \right\|_{\HS}^2,
 \label{G:eq:completeleak}
\end{equation}
which is the complete-shell leakage scalar of
Ref.~\cite{BateninPreece2026}.
\end{corollary}

\begin{proof}
At such a vertex, $P_{\mathbf L}$ is the orthogonal projector onto the union
of the selected shells.  Orthogonality of the signed-shell decomposition
splits the squared leakage norm into the sum in \eqref{G:eq:completeleak}.
\end{proof}

\begin{remark}
The complete-shell criterion tests finitely many distinguished vertices of a
much larger moduli problem.  The function \eqref{G:eq:leakage} extends that
test to every partial-shell configuration with the same finite spectral
support.
\end{remark}

\subsection{Infinitesimal deformation of the leakage variety}

The algebraic zero-locus description also has a concrete first-order
deformation theory.  We record it in a graph chart, which avoids assigning a
spurious codimension to a singular zero set.

Let
\[
 W=L_\alpha\oplus L_\beta
 \subset E_\alpha\oplus E_\beta
\]
be a zero-leakage point for two distinct signed shells.  A tangent vector to
the product Grassmannian is represented by
\[
 X_\alpha\in\operatorname{Hom}(L_\alpha,L_\alpha^\perp),
 \qquad
 X_\beta\in\operatorname{Hom}(L_\beta,L_\beta^\perp).
\]

\begin{theorem}[Linearized leakage and Zariski tangent space]
\label{G:thm:linearizedleakage}
In the standard graph trivialization of the leakage bundle, the differential
at a zero-leakage point is
\begin{align}
 d\Phi(X_\alpha,X_\beta)(u,v)
 ={}&
 (I-P_W)
 \left[
 Q(X_\alpha u,v)+Q(u,X_\beta v)
 \right]
 \notag\\
 &-
 X_\alpha P_{L_\alpha}Q(u,v)
 -
 X_\beta P_{L_\beta}Q(u,v),
 \label{G:eq:linearizedleakage}
\end{align}
for \(u\in L_\alpha\), \(v\in L_\beta\).  Consequently
\begin{equation}
 \boxed{
 T_{(L_\alpha,L_\beta)}^{\rm Zar}\mathcal Z_{\alpha,\beta}
 =
 \ker d\Phi.
 }
 \label{G:eq:zartangent}
\end{equation}
At a regular point the local codimension equals
\(\operatorname{rank}d\Phi\).
\end{theorem}

\begin{proof}
Represent nearby planes as graphs of
\(\varepsilon X_\alpha\) and \(\varepsilon X_\beta\).
Differentiate the bilinear leakage
\[
 (I-P_{W(\varepsilon)})
 Q(u+\varepsilon X_\alpha u,
   v+\varepsilon X_\beta v)
\]
at \(\varepsilon=0\).  The first line of
\eqref{G:eq:linearizedleakage} comes from varying the two arguments, while the
second line is the first variation of the target projection.  The tangent
space to the zero locus of an algebraic section is the kernel of its
linearization.  At a regular zero, the constant-rank theorem gives the
codimension statement.
\end{proof}

\begin{remark}
There is no universal formula equating the codimension of the leakage variety
with the sum of ranks of the shell projections of \(Q\).  Such a formula
would require additional transversality and independence hypotheses.
\end{remark}

\subsection{Pressure-linear incidence varieties}

One may impose the stronger requirement
\[
          Q(L,L)=0
\]
rather than merely $Q(L,L)\subset L$.  For two distinct shells
$E_\alpha,E_\beta$, same-shell terms already vanish.  Hence pressure-linearity
of $L_\alpha\oplus L_\beta$ is equivalent to
\[
          Q(L_\alpha,L_\beta)=0.
\]

\begin{definition}
For dimensions $d,e$, define
\begin{equation}
 \mathcal I_{\alpha,\beta}^{d,e}
 :=
 \{(L,M)\in
 \Gr(d,E_\alpha)\times\Gr(e,E_\beta):
 Q(L,M)=0\}.
 \label{G:eq:incidence}
\end{equation}
\end{definition}

\begin{proposition}[Cross-shell incidence variety]
\label{G:prop:incidence}
The set $\mathcal I_{\alpha,\beta}^{d,e}$ is a compact
$SO(4)$-invariant real-algebraic subset of the product Grassmannian.
Every point parameterizes a pressure-linear exact reduction
$L\oplus M$.
\end{proposition}

\begin{proof}
Use projector coordinates as in \cref{G:thm:leakagevariety}.  The condition
$Q(L,M)=0$ is equivalent to the vanishing of the finite family of bilinear
coefficients
\[
      Q(P_Le_i,P_Mf_j),
\]
hence is algebraic.  Stokes invariance is automatic because $L$ and $M$ lie
in eigenspaces of $A$; the same-shell terms vanish by
\cref{G:lem:shellcancel}.
\end{proof}

This incidence geometry separates two questions.  Representation theory
determines which shell pairs can interact at all; algebraic geometry then
determines which partial subspaces annihilate the allowed interaction tensor.


\section{Killing precession on reduction moduli}

\subsection{Killing drift on a signed shell}

Let \(E_{k,\sigma}\), \(k\ge1\), \(\sigma=\pm\), be a real signed curl
eigenspace on the round three-sphere of radius \(R\).  Its dimension is
\[
 d_k=(k+1)(k+3),
\]
and the Stokes operator is scalar:
\[
 A|_{E_{k,\sigma}}=\kappa_k\Id,
 \qquad
 \kappa_k=\frac{k(k+4)}{R^2}.
\]
The full Killing algebra splits as
\[
 \cK=\cK_+\oplus\cK_-
 \simeq\mathfrak{su}(2)_+\oplus\mathfrak{su}(2)_-.
\]

We use the exact shell drift
\begin{equation}
\begin{aligned}
 D_{k,\sigma}(m)
 &=
 a_k\,\ad_{m_\sigma}
 +
 b_k\,\ad_{m_{-\sigma}},\\
 a_k&=\frac{k}{k+2},
 \qquad
 b_k=\frac{k+4}{k+2}.
\end{aligned}
\label{G:pr:eq:D}
\end{equation}
Here \(\ad\) denotes the infinitesimal isometry action on the shell.
The two simple factors commute, and \(D_{k,\sigma}(m)\) is real
skew-adjoint.

The complexification of \(E_{k,\sigma}\) is the tensor product
\[
 V_{j_\sigma}\boxtimes V_{j_{-\sigma}},
 \qquad
 j_\sigma=\frac{k+2}{2},
 \quad
 j_{-\sigma}=\frac{k}{2}.
\]
These are the only representation-theoretic inputs used below.

\subsection{Projector Lax dynamics}

Let \(L_0\subset E_{k,\sigma}\) be a real \(d\)-dimensional subspace and let
\(U(t)\) solve
\[
 \dot U(t)=-D(t)U(t),
 \qquad U(0)=I,
\]
where \(D(t)=D_{k,\sigma}(m(t))\).  Since \(D(t)^*=-D(t)\), \(U(t)\) is
orthogonal.

Define
\[
 L(t)=U(t)L_0,
 \qquad
 P(t)=U(t)P_0U(t)^*.
\]

\begin{theorem}[Exact projector equation]
\label{G:pr:thm:lax}
The orthogonal projector satisfies
\begin{equation}
 \boxed{\dot P=[P,D].}
 \label{G:pr:eq:lax}
\end{equation}
If \(D\) is constant, then
\begin{equation}
 L_D(t):=P(t)D
\end{equation}
satisfies the Lax equation
\begin{equation}
 \boxed{\dot L_D=[L_D,D].}
 \label{G:pr:eq:LD}
\end{equation}
Consequently
\begin{equation}
 \tr\bigl(L_D(t)^r\bigr)
 =
 \tr\bigl(L_D(0)^r\bigr)
 \qquad(r\ge1).
 \label{G:pr:eq:invariants}
\end{equation}
\end{theorem}

\begin{proof}
Differentiate \(P=UP_0U^*\):
\[
 \dot P=-DP+PD=[P,D].
\]
For constant \(D\),
\[
 \dot L_D=[P,D]D=PD^2-DPD=[PD,D].
\]
Trace invariance follows from cyclicity.
\end{proof}

\begin{remark}
The invariants \eqref{G:pr:eq:invariants} are spectral invariants of one Lax
matrix.  Their existence does not by itself establish Liouville or
Gelfand--Tsetlin complete integrability.
\end{remark}

\subsection{Exact geodesic criterion}

The physical space of real \(d\)-planes is
\[
 \Gr_{\mathbb R}(d,d_k)
 \simeq
 O(d_k)/(O(d)\times O(d_k-d)),
\]
with its canonical metric.  Geodesics in the projector model are generated
by horizontal, block-off-diagonal skew matrices
\cite{Edelman1998,Bendokat2024}.

Assume now that \(D\) is constant.  Relative to the orthogonal splitting
\[
 E_{k,\sigma}=L_0\oplus L_0^\perp,
\]
write
\begin{equation}
 D=
 \begin{pmatrix}
 A&-B^*\\
 B&C
 \end{pmatrix},
 \qquad
 A^*=-A,\quad C^*=-C.
 \label{G:pr:eq:blockD}
\end{equation}

\begin{theorem}[Geodesic criterion]
\label{G:pr:thm:geodesic}
The orbit
\[
 P(t)=e^{-tD}P_0e^{tD}
\]
is a canonical Grassmann geodesic if and only if
\begin{equation}
 \boxed{CB=BA.}
 \label{G:pr:eq:geodesiccriterion}
\end{equation}
Equivalently, the block-diagonal and block-off-diagonal parts of \(D\)
commute.
\end{theorem}

\begin{proof}
At \(P_0=\mathrm{diag}(I,0)\),
\[
 \dot P(0)=
 \begin{pmatrix}
 0&-B^*\\
 -B&0
 \end{pmatrix}.
\]
A second differentiation gives
\[
 \ddot P(0)=
 \begin{pmatrix}
 -2B^*B&A B^*-B^*C\\
 CB-BA&2BB^*
 \end{pmatrix}.
\]
For the canonical Grassmann metric, the tangential part of an ambient
symmetric matrix is precisely its block-off-diagonal part.  Thus vanishing
covariant acceleration at \(t=0\) is equivalent to
\(CB=BA\).

If this condition holds, the decomposition
\(D=D_{\rm diag}+D_{\rm off}\) satisfies
\[
 [D_{\rm diag},D_{\rm off}]=0,
 \qquad
 [D_{\rm diag},P_0]=0.
\]
Hence
\[
 e^{-tD}P_0e^{tD}
 =
 e^{-tD_{\rm off}}P_0e^{tD_{\rm off}},
\]
which is the standard Grassmann geodesic generated by the horizontal
skew matrix \(D_{\rm off}\).  This proves sufficiency as well.
\end{proof}

\begin{corollary}
A generic constant Killing drift orbit is not a Grassmann geodesic.
Geodesicity is an algebraic compatibility condition between the action
inside \(L_0\), the action inside \(L_0^\perp\), and the mixing block \(B\).
\end{corollary}

\subsection{Exact fixed-drift spectrum}

Choose pointwise-orthonormal Cartan generators
\(X_3^\pm\) satisfying the standard round-sphere normalization.  By
independent rotations in the two simple factors, any fixed Killing momentum
may be written
\[
 m_\sigma=M_\sigma X_3^\sigma,
 \qquad
 m_{-\sigma}=M_{-\sigma}X_3^{-\sigma},
 \qquad
 M_\pm\ge0.
\]

On the complexified shell choose the tensor-product weight basis
\[
 |p,q\rangle,
 \qquad
 p=-j_\sigma,-j_\sigma+1,\ldots,j_\sigma,
\]
\[
 q=-j_{-\sigma},-j_{-\sigma}+1,\ldots,j_{-\sigma}.
\]

\begin{theorem}[Exact shell splitting]
\label{G:pr:thm:spectrum}
The operator
\[
 \mathcal L_m=\nu A+D_{k,\sigma}(m)
\]
has eigenvalues
\begin{equation}
 \boxed{
 \lambda_{p,q}
 =
 \nu\kappa_k
 +
 \frac{2i}{R}
 \left(
 a_kM_\sigma p+
 b_kM_{-\sigma}q
 \right).
 }
 \label{G:pr:eq:spectrum}
\end{equation}
\end{theorem}

\begin{proof}
The two \(SU(2)\) factors commute.  The infinitesimal generator
\(X_3^\tau\) has eigenvalues \(2ip/R\) in spin \(j_\tau\).  Apply
\eqref{G:pr:eq:D} and add the scalar Stokes part.
\end{proof}

\subsection{Resonance arithmetic and generic simplicity}

Two weights \((p,q)\) and \((p',q')\) collide exactly when
\[
 a_kM_\sigma(p-p')
 +
 b_kM_{-\sigma}(q-q')
 =
 0.
\]
The differences are integers:
\[
 r=p-p',
 \qquad
 s=q-q',
\]
with
\[
 |r|\le k+2,
 \qquad
 |s|\le k.
\]

\begin{theorem}[Finite Diophantine collision criterion]
\label{G:pr:thm:collision}
The spectrum \eqref{G:pr:eq:spectrum} is simple if and only if
\begin{equation}
 a_kM_\sigma r+b_kM_{-\sigma}s\ne0
 \label{G:pr:eq:nocollision}
\end{equation}
for every nonzero admissible integer pair
\[
 (r,s)\in\mathbb Z^2,
 \qquad
 |r|\le k+2,\quad |s|\le k.
\]
If both \(M_\sigma\) and \(M_{-\sigma}\) are nonzero, an irrational ratio
\[
 \frac{a_kM_\sigma}{b_kM_{-\sigma}}
\]
is sufficient for simplicity.
\end{theorem}

\begin{corollary}[Exact minimum spacing]
\label{G:pr:cor:spacing}
The minimum nonzero frequency spacing is
\begin{equation}
 \boxed{
 \Delta\omega_k
 =
 \frac2R
 \min_{\substack{(r,s)\in\mathbb Z^2\setminus\{0\}\\
 |r|\le k+2,\ |s|\le k}}
 |a_kM_\sigma r+b_kM_{-\sigma}s|.
 }
 \label{G:pr:eq:spacing}
\end{equation}
Thus the minimum spacing can be smaller than either one-axis step because of
cross-factor cancellation.
\end{corollary}

\begin{remark}
Spectral splitting alone does not prove suppression of nonlinear triads.
A cascade estimate requires a separate oscillatory-integral theorem and
explicit exclusion of resonances satisfying \eqref{G:pr:eq:nocollision} with
equality.
\end{remark}

\subsection{Special-orthogonal shell transport}

The shell drift is a real skew-adjoint representation operator.

\begin{lemma}[Tracelessness]
\label{G:pr:lem:traceless}
For every \(m\in\cK\),
\begin{equation}
 \boxed{\tr D_{k,\sigma}(m)=0.}
 \label{G:pr:eq:traceless}
\end{equation}
\end{lemma}

\begin{proof}
Each simple-factor infinitesimal representation of
\(\mathfrak{su}(2)\) is traceless.  The drift is a linear combination of the
two commuting simple-factor actions.
\end{proof}

\begin{theorem}[Determinant-one transport]
\label{G:pr:thm:detone}
Let
\[
 \dot U(t)=-D_{k,\sigma}(m(t))U(t),
 \qquad U(0)=I.
\]
Then
\begin{equation}
 \boxed{
 U(t)\in SO(d_k),
 \qquad
 \det U(t)=1
 }
 \label{G:pr:eq:detone}
\end{equation}
for every \(t\).  After complexification,
\(U(t)\in SU(d_k)\).
\end{theorem}

\begin{proof}
Skew-adjointness gives orthogonality.  Jacobi's determinant formula and
\cref{G:pr:lem:traceless} give
\[
 \frac{\dd}{\dd t}\log\det U
 =
 -\tr D_{k,\sigma}(m(t))
 =
 0.
\]
Since \(\det U(0)=1\), the result follows.
\end{proof}

\begin{corollary}[No determinant Berry phase]
\label{G:pr:cor:nouone}
The determinant line of the full shell has trivial parallel transport.
Therefore full-shell precession carries no nontrivial scalar \(U(1)\)
Berry phase obtained from \(\det U\), even though the nonabelian
\(SO(d_k)\) holonomy can be nontrivial.
\end{corollary}

\begin{remark}[Closed loops]
A periodic coefficient \(m(T)=m(0)\) does not by itself define holonomy of a
closed base loop.  If the connection is viewed over \(SO(4)\), the
reconstructed group path must satisfy \(g(T)=g(0)\).  Likewise, to speak of
a holonomy operator on a particular subspace \(L_0\), the endpoint transport
must return that subspace to itself.
\end{remark}

\subsection{Curvature without scalar topology}

The shell curvature may be written
\[
\begin{aligned}
 \Omega_{k,\sigma}(m,n)
 &=
 -\frac{2k}{(k+2)^2}
 \ad_{[m_\sigma,n_\sigma]}\\
 &\quad+
 \frac{2(k+4)}{(k+2)^2}
 \ad_{[m_{-\sigma},n_{-\sigma}]}.
\end{aligned}
\]
Every term is traceless, hence
\[
 \tr\Omega_{k,\sigma}=0.
\]
Thus nonflatness is a genuinely nonabelian phenomenon.  It should not be
confused with a nonzero first Chern class or with a quantized scalar
Chern--Simons charge.  Chern--Simons functionals require their own
normalization and gauge analysis; in general their real values are not
topological integers.


\section{Killing-response transform and minimal observability}

\subsection{Killing--Beltrami response identity}

Choose pointwise-orthonormal frames $X_a^\tau$ in the two simple Killing
factors with
\[
 [X_a^\tau,X_b^\tau]
 =
 \frac{2}{R}\varepsilon_{abc}X_c^\tau,
 \qquad
 [X_a^+,X_b^-]=0,
\]
and put $K_a^\tau=V_R^{-1/2}X_a^\tau$.

\begin{lemma}[Killing--Beltrami response]
\label{R:lem:KB}
If $K\in E_{0,\tau}$ and $v\in E_{k,\sigma}$, $k\ge1$, then
\begin{equation}
 \boxed{
 2Q(K,v)=
 \left(1-\frac{2\tau\sigma}{k+2}\right)[K,v].
 }
 \label{R:eq:KB}
\end{equation}
In particular $Q(K,v)\in E_{k,\sigma}$.
\end{lemma}

\begin{proof}
On an oriented Riemannian three-manifold,
\[
 \nabla_Kv+\nabla_vK
 =
 \gradg\langle K,v\rangle
 -K\times\curl v-v\times\curl K.
\]
With
$\lambda=2\tau/R$ and
$\mu=\sigma(k+2)/R$,
the non-gradient part is
$(\lambda-\mu)K\times v$.
Since $K$ is Killing, its flow commutes with curl, so
$[K,v]\in E_{k,\sigma}$.  Also
$\curl(K\times v)=-[K,v]$.
After Leray projection,
\[
 \PP(K\times v)=-\mu^{-1}[K,v].
\]
Substitution gives \eqref{R:eq:KB}.
\end{proof}

The representation of the two $SU(2)$ factors on $E_{k,\sigma}$ has spins
\begin{equation}
       j_\tau=\frac{k+1+\tau\sigma}{2}.
       \label{R:eq:spin}
\end{equation}
For the $L^2$-normalized Killing basis,
\begin{equation}
 -\sum_{a=1}^3\mathcal L_{K_a^\tau}^2
 =
 \frac{4}{R^2V_R}j_\tau(j_\tau+1)\Id.
 \label{R:eq:partialCasimir}
\end{equation}

\subsection{The six-channel Parseval frame}

Set
\[
       T_a^\tau=Q(K_a^\tau,\cdot).
\]
These operators are skew-adjoint and preserve every signed shell.

Define
\[
 (\cA v)_{a,\tau}
 =
 \sqrt{2V_R}\,T_a^\tau v
\]
on $D(A^{1/2})\cap\mathcal H$.

\begin{theorem}[Casimir factorization]
\label{R:thm:factor}
On $\mathcal H$,
\begin{equation}
 \boxed{\cA^*\cA=A.}
 \label{R:eq:factor}
\end{equation}
Equivalently,
\begin{equation}
 \boxed{
 2V_R\sum_{\tau,a}\|T_a^\tau v\|_2^2
 =
 \langle Av,v\rangle.
 }
 \label{R:eq:diss}
\end{equation}
\end{theorem}

\begin{proof}
On each signed shell, use \eqref{R:eq:KB} and
\eqref{R:eq:partialCasimir}.  The sum of the two chiral scalar contributions
reduces to $\kappa_k$.  Shell orthogonality then gives the operator identity.
\end{proof}


The same identity may be written in sum-of-squares form as
\begin{equation}
 \boxed{
 A
 =
 -2V_R
 \sum_{\tau=\pm}\sum_{a=1}^3
 (T_a^\tau)^2.
 }
 \label{R:eq:CasimirCore}
\end{equation}

Define the bounded normalized transforms
\begin{equation}
 \boxed{
 \cR_a^\tau
 =
 \sqrt{2V_R}\,T_a^\tau A^{-1/2}
 :
 \mathcal H\to L^2_\sigma(T\Sthree).
 }
 \label{R:eq:R}
\end{equation}

\begin{theorem}[Parseval Killing-response frame]
\label{R:thm:parseval}
\begin{equation}
 \boxed{
 \sum_{\tau=\pm}\sum_{a=1}^3
 (\cR_a^\tau)^*\cR_a^\tau
 =
 I_{\mathcal H}.
 }
 \label{R:eq:parseval}
\end{equation}
Hence
\[
       \|v\|_2^2
       =
       \sum_{\tau,a}\|\cR_a^\tau v\|_2^2.
\]
\end{theorem}

\begin{proof}
Conjugate \eqref{R:eq:factor} by $A^{-1/2}$.
\end{proof}

\subsection{A single chirality is already complete}

Let
\[
       \cA_\tau v
       =
       \sqrt{2V_R}
       (T_1^\tau v,T_2^\tau v,T_3^\tau v).
\]

\begin{theorem}[Exact chiral frame operator]
\label{R:thm:chiral}
On $\mathcal H$,
\begin{align}
 \boxed{
 \cA_\tau^*\cA_\tau
 =
 \frac12A
 \left(I-\frac{2\tau}{R}\curl^{-1}\right),
 }
 \label{R:eq:chiralraw}\\
 \boxed{
 \cS_\tau:=
 \sum_{a=1}^3(\cR_a^\tau)^*\cR_a^\tau
 =
 \frac12
 \left(I-\frac{2\tau}{R}\curl^{-1}\right).
 }
 \label{R:eq:chiralnorm}
\end{align}
Moreover,
\begin{equation}
 \boxed{
 \frac16I\le\cS_\tau\le\frac56I,
 }
 \label{R:eq:sharpbounds}
\end{equation}
and both constants are sharp.
\end{theorem}

\begin{proof}
On $E_{k,\sigma}$, \eqref{R:eq:KB} and
\eqref{R:eq:partialCasimir} give
\[
 \cA_\tau^*\cA_\tau
 =
 \frac{\kappa_k}{2}
 \left(1-\frac{2\tau\sigma}{k+2}\right)\Id.
\]
Since
$\curl^{-1}=\sigma R/(k+2)$ on this shell,
\eqref{R:eq:chiralraw} follows.  Conjugation by $A^{-1/2}$ gives
\eqref{R:eq:chiralnorm}.  Finally,
\[
       \frac{2}{k+2}\le\frac23,
       \qquad k\ge1,
\]
so the spectral interval is $[1/6,5/6]$.  Both endpoints occur at $k=1$.
\end{proof}

\begin{corollary}[Stable three-channel reconstruction]
\label{R:cor:three}
For either chirality,
\begin{equation}
 \frac16\|v\|_2^2
 \le
 \sum_{a=1}^3\|\cR_a^\tau v\|_2^2
 \le
 \frac56\|v\|_2^2.
 \label{R:eq:three}
\end{equation}
The canonical dual reconstruction is
\begin{equation}
 \boxed{
 v
 =
 \cS_\tau^{-1}
 \sum_{a=1}^3
 (\cR_a^\tau)^*\cR_a^\tau v.
 }
 \label{R:eq:dual}
\end{equation}
\end{corollary}

For raw response fields,
\begin{equation}
 \boxed{
 v
 =
 -4V_R
 \left[
 A\left(I-\frac{2\tau}{R}\curl^{-1}\right)
 \right]^{-1}
 \sum_{a=1}^3
 T_a^\tau(T_a^\tau v).
 }
 \label{R:eq:rawreconstruction}
\end{equation}

\subsection{Two probes determine the entire high-mode field}

Fix an orthonormal pair
$K_1^\tau,K_2^\tau\in\cK_\tau$ and define
\[
       \cA_{\tau,2}v
       =
       \sqrt{2V_R}(T_1^\tau v,T_2^\tau v).
\]

\begin{lemma}[Two-axis shell spectrum]
\label{R:lem:pair}
On $E_{k,\sigma}$ set
\[
       j=\frac{k+1+\tau\sigma}{2},
\qquad
       d_\tau=
       \frac{\kappa_k}{2}
       \left(1-\frac{2\tau\sigma}{k+2}\right).
\]
In a weight basis of the omitted third generator, the pair frame eigenvalues
are
\begin{equation}
 \lambda_m^{(2)}
 =
 d_\tau
 \frac{j(j+1)-m^2}{j(j+1)},
 \qquad
 m=-j,\ldots,j.
 \label{R:eq:pairlambda}
\end{equation}
Hence
\begin{equation}
 \boxed{
 \lambda_{\min}^{(2)}
 =
 \frac{d_\tau}{j+1}>0.
 }
 \label{R:eq:pairmin}
\end{equation}
\end{lemma}

\begin{proof}
The common scalar in \eqref{R:eq:KB} cancels in the ratio between the full
three-axis Casimir and the two-axis operator.  In spin $j$,
\[
       J_1^2+J_2^2=j(j+1)I-J_3^2.
\]
The minimum occurs at $|m|=j$, where the numerator is $j$.
\end{proof}

\begin{theorem}[Sharp global two-probe bound]
\label{R:thm:two}
For either chirality and every
$v\in D(A^{1/2})\cap\mathcal H$,
\begin{equation}
 \boxed{
 2V_R
 \left(
 \|T_1^\tau v\|_2^2+
 \|T_2^\tau v\|_2^2
 \right)
 \ge
 \frac{1}{3R^2}\|v\|_2^2.
 }
 \label{R:eq:twobound}
\end{equation}
The constant is sharp.  Thus two fixed same-chirality Killing responses
determine the complete non-Killing state and have a bounded inverse on their
range in the $L^2$ topology.
\end{theorem}

\begin{proof}
By shell orthogonality, minimize \eqref{R:eq:pairmin}.

If $\tau\sigma=+1$, then
\[
 j=\frac{k+2}{2},
 \qquad
 \lambda_{\min}^{(2)}
 =
 \frac{k^2}{R^2(k+2)}.
 \label{R:eq:same}
\]
This is increasing for $k\ge1$ and equals $1/(3R^2)$ at $k=1$.

If $\tau\sigma=-1$, then
\[
 j=\frac{k}{2},
 \qquad
 \lambda_{\min}^{(2)}
 =
 \frac{k(k+4)^2}{R^2(k+2)^2},
 \label{R:eq:opp}
\]
which is larger than $1/(3R^2)$ for all $k\ge1$.
\end{proof}

\begin{remark}
The raw pair is a first-order closed analysis operator, not a bounded map
$L^2\to(L^2)^2$.  The estimate \eqref{R:eq:twobound} states that its inverse
from the exact response range to $L^2$ is bounded.  The normalized
three- and six-channel transforms are the bounded $g$-frames.
\end{remark}

\subsection{One-probe parity law}

Fix $0\ne K\in\cK_\tau$.

\begin{theorem}[Parity and exact blind-space dimension]
\label{R:thm:parity}
On $E_{k,\sigma}$:

\begin{enumerate}[label=(\roman*)]
\item If $k$ is odd, then
\[
       T_K=Q(K,\cdot)
\]
is invertible.

\item If $k$ is even, then $T_K$ has a nontrivial kernel of dimension
\begin{equation}
 \boxed{
 \dim\ker T_K
 =
 k+2-\tau\sigma.
 }
 \label{R:eq:kerneldim}
\end{equation}
Thus the dimensions are $k+1$ for $\tau\sigma=+1$ and $k+3$ for
$\tau\sigma=-1$.
\end{enumerate}
\end{theorem}

\begin{proof}
By \eqref{R:eq:KB}, $T_K$ is a nonzero scalar multiple of an
$SU(2)_\tau$ infinitesimal generator in spin
\[
       j_\tau=\frac{k+1+\tau\sigma}{2}.
\]
After complexification its weights are
$-j_\tau,-j_\tau+1,\ldots,j_\tau$.
Zero occurs if and only if $j_\tau$ is an integer.

For odd $k$, both possible spins are half-integers.
For even $k$, both are integers.  The zero-weight line in the
$SU(2)_\tau$ factor is tensored with the opposite factor of dimension
\[
       2j_{-\tau}+1
       =
       k+2-\tau\sigma.
\]
Complexification preserves the kernel dimension of the real operator.
\end{proof}

\begin{corollary}[One-probe condition number on odd shells]
\label{R:cor:onecond}
For odd $k$,
\begin{equation}
 \boxed{
 \operatorname{cond}(T_K)
 =
 k+1+\tau\sigma.
 }
 \label{R:eq:onecond}
\end{equation}
Thus the condition number is $k+2$ for $\tau\sigma=+1$ and $k$ for
$\tau\sigma=-1$.
\end{corollary}

\begin{proof}
For half-integer spin $j$, the largest generator weight has magnitude $j$
and the smallest has magnitude $1/2$.  The ratio is $2j$.
\end{proof}

\subsection{Two-probe condition numbers and minimality}

\begin{theorem}[Exact pair condition number]
\label{R:thm:paircond}
For two orthonormal same-chirality probes on $E_{k,\sigma}$, with
$j=(k+1+\tau\sigma)/2$,
\begin{equation}
 \boxed{
 \operatorname{cond}_2
 =
 \begin{cases}
 \sqrt{j+1},&j\in\mathbb Z,\\[1ex]
 \sqrt{j+1-\dfrac{1}{4j}},
 &j\in\mathbb Z+\dfrac12.
 \end{cases}
 }
 \label{R:eq:paircond}
\end{equation}
Hence the pair condition number is $O(k^{1/2})$, whereas one-channel
inversion on odd shells is $O(k)$.
\end{theorem}

\begin{proof}
The squared singular values are proportional to
$j(j+1)-m^2$ by \eqref{R:eq:pairlambda}.  The minimum is $j$.
For integer $j$, the maximum is $j(j+1)$ at $m=0$.
For half-integer $j$, it is $j(j+1)-1/4$ at $|m|=1/2$.
\end{proof}

\begin{corollary}[Minimal universal probe count]
\label{R:cor:minimal}
Within either fixed chiral Killing factor:

\begin{enumerate}[label=(\roman*)]
\item two linearly independent fixed probes are sufficient to distinguish
every non-Killing state;

\item one fixed probe is never sufficient uniformly over all shells because
every even shell has a blind subspace;

\item on a prescribed odd shell the minimal number is one, while on a
prescribed even shell it is two.
\end{enumerate}
\end{corollary}

\subsection{Erasure hierarchy}

The reconstruction structure can be summarized as follows.

\begin{theorem}[Chiral erasure hierarchy]
\label{R:thm:erasure}
Fix $E_{k,\sigma}$ and one Killing chirality.

\begin{enumerate}[label=(\roman*)]
\item Three orthonormal response channels form a tight operator frame on the
shell.

\item Erasing any one channel leaves a frame with lower eigenvalue
$d_\tau/(j_\tau+1)$.

\item Erasing two channels leaves an injective one-channel transform if and
only if $k$ is odd.

\item On even shells, every one-channel transform has the blind-space
dimension \eqref{R:eq:kerneldim}.
\end{enumerate}
\end{theorem}

The result gives exact loss tolerance of the Killing-response representation.


\section{Casimir response geometry}

\subsection{From lossless response fields to Gram compression}

The frame results above use the response fields themselves.  We now compress
those six fields to their mutual inner products.  This passage is useful but
not lossless.  It produces finite-dimensional isometry covariants and exact
spectral statistics, while forgetting part of the common phase information.
This hierarchy should be kept explicit:
\[
\boxed{
v
\longmapsto
\{Q(K_a^\tau,v)\}_{\tau,a}
\longmapsto
\mathscr M_v
\longmapsto
\{\mathscr M_v^{[s]}\}_{s}
\longmapsto
\Phi_v(\theta).
}
\]
The first arrow is injective by the two-probe theorem.  Later we show that the
Gram-level map cannot be injective on sufficiently high individual shells.

\subsection{The Casimir response tensor}

\begin{definition}[Response operator and response tensor]
\label{R:rg:def:response}
For a divergence-free field $u$ in the form domain of $A$, define
\begin{equation}
 \cB_u:\mathfrak K\longrightarrow L^2_\sigma(T\Sthree),
 \qquad
 \cB_uX
 :=
 \sqrt{2V_R}\,Q(X,u).
 \label{R:rg:eq:Rdef}
\end{equation}
Its Casimir response tensor is
\begin{equation}
 \boxed{
 \cM_u
 :=
 \cB_u^*\cB_u
 }
 \label{R:rg:eq:Moperator}
\end{equation}
or, equivalently,
\begin{equation}
 \cM_u(X,Y)
 =
 2V_R
 \langle Q(X,u),Q(Y,u)\rangle_{L^2}.
 \label{R:rg:eq:Mform}
\end{equation}
\end{definition}

The tensor is positive semidefinite and quadratic in $u$:
\[
      \cM_{cu}=c^2\cM_u.
\]

\begin{theorem}[Trace identity and exact dissipation]
\label{R:rg:thm:trace}
For every divergence-free $u$ in the form domain of $A$,
\begin{equation}
 \boxed{
 \tr_{\mathfrak K}\cM_u
 =
 \langle Au,u\rangle
 =
 2\|\Def u\|_2^2.
 }
 \label{R:rg:eq:trace}
\end{equation}
Hence the viscous energy dissipation satisfies
\begin{equation}
 \boxed{
 \varepsilon_\nu(u)
 :=
 2\nu\|\Def u\|_2^2
 =
 \nu\,\tr\cM_u.
 }
 \label{R:rg:eq:diss}
\end{equation}
\end{theorem}

\begin{proof}
The quadratic Casimir identity for the six $L^2$-normalized Killing
responses is
\begin{equation}
 A
 =
 -2V_R
 \sum_{\tau=\pm}\sum_{a=1}^3
 T_a^\tau T_a^\tau,
 \qquad
 T_a^\tau:=Q(K_a^\tau,\cdot),
 \label{R:rg:eq:Acasimir}
\end{equation}
where each $T_a^\tau$ is skew-adjoint.  Therefore
\[
 \langle Au,u\rangle
 =
 2V_R
 \sum_{\tau,a}
 \|Q(K_a^\tau,u)\|_2^2
 =
 \tr\cM_u.
\]
For the deformation-tensor viscosity,
$\langle Au,u\rangle=2\|\Def u\|_2^2$.
\end{proof}

\begin{remark}
The response tensor is a global $L^2$ object.  Equation
\eqref{R:rg:eq:trace} does not, by itself, imply spatial localization,
microlocal anisotropy, or a preferred geometric axis in physical space.
\end{remark}

\subsection{Isometry covariance and the response ellipsoid}

Let $g\in SO(4)$ act unitarily on vector fields by pushforward.  The induced
action on $\mathfrak K$ is the adjoint action.

\begin{theorem}[Isometry covariance]
\label{R:rg:thm:covariance}
For every $g\in SO(4)$,
\begin{equation}
 \cM_{g_*u}
 =
 \operatorname{Ad}_g\,
 \cM_u\,
 \operatorname{Ad}_{g^{-1}}.
 \label{R:rg:eq:covariance}
\end{equation}
Consequently the ordered eigenvalues
\begin{equation}
 \lambda_1(u)\ge\cdots\ge\lambda_6(u)\ge0
 \label{R:rg:eq:eigenvalues}
\end{equation}
are invariants of the physical isometry orbit of $u$.
\end{theorem}

\begin{proof}
Isometry equivariance of the Levi-Civita connection and Leray projection
gives
\[
 Q(X,g_*u)
 =
 g_*Q(\operatorname{Ad}_{g^{-1}}X,u).
\]
Since $g_*$ is unitary in $L^2$, substitution into \eqref{R:rg:eq:Mform} gives
\eqref{R:rg:eq:covariance}.
\end{proof}

\begin{definition}[Response ellipsoid]
The image of the unit sphere in $\mathfrak K$ under $\cB_u$ is the response
ellipsoid in $\operatorname{ran}\cB_u$.  Its principal semiaxis lengths are
\begin{equation}
        \sqrt{\lambda_1(u)},\ldots,\sqrt{\lambda_6(u)},
 \label{R:rg:eq:axes}
\end{equation}
with zero axes allowed.
\end{definition}

\begin{remark}
The eigenvalues themselves are \emph{squared} semiaxis lengths.  This
distinction is important when interpreting the geometry.
\end{remark}

\begin{corollary}[Extremal-response principle]
\label{R:rg:cor:extremal}
For every $u$,
\begin{equation}
 \frac{1}{6}\langle Au,u\rangle
 \le
 \lambda_1(u)
 \le
 \langle Au,u\rangle.
 \label{R:rg:eq:lambdabounds}
\end{equation}
Hence along any sequence $u_n$,
\begin{equation}
 \|\Def u_n\|_2\longrightarrow\infty
 \quad\Longleftrightarrow\quad
 \lambda_1(u_n)\longrightarrow\infty.
 \label{R:rg:eq:blowupequiv}
\end{equation}
\end{corollary}

\begin{proof}
For a positive semidefinite $6\times6$ operator,
\[
       \lambda_1
       \le\sum_i\lambda_i
       \le6\lambda_1.
\]
Use \cref{R:rg:thm:trace}.
\end{proof}

\begin{remark}
Equation \eqref{R:rg:eq:blowupequiv} does not say that dissipation growth must
become anisotropic.  All six eigenvalues may grow simultaneously.  A
separate normalized spectral diagnostic is required to discuss response
polarization.
\end{remark}

\subsection{Spectral additivity}

Write
\[
       u=\sum_{k,\sigma}u_{k,\sigma},
       \qquad
       u_{k,\sigma}\in E_{k,\sigma}.
\]

\begin{theorem}[Spectral additivity of the response tensor]
\label{R:rg:thm:additivity}
For every $u$ in the form domain of $A$,
\begin{equation}
 \boxed{
 \cM_u
 =
 \sum_{k,\sigma}
 \cM_{u_{k,\sigma}}
 }
 \label{R:rg:eq:additivity}
\end{equation}
with convergence in the finite-dimensional operator norm on $\mathfrak K$
whenever $\langle Au,u\rangle<\infty$.
\end{theorem}

\begin{proof}
By \eqref{R:eq:KB},
\[
 Q(X,u_{k,\sigma})\in E_{k,\sigma}
\]
for every $X\in\mathfrak K$.  Distinct signed shells are $L^2$-orthogonal.
Hence for $X,Y\in\mathfrak K$,
\[
 \langle Q(X,u),Q(Y,u)\rangle
 =
 \sum_{k,\sigma}
 \langle
 Q(X,u_{k,\sigma}),
 Q(Y,u_{k,\sigma})
 \rangle.
\]
This is \eqref{R:rg:eq:additivity}.  The trace of the positive series is
$\langle Au,u\rangle$, so convergence follows in finite dimension.
\end{proof}

This identity says that shell phases can affect the shape of each
$\cM_{u_{k,\sigma}}$, but there are no cross-shell interference terms in
the response tensor itself.

\subsection{Exact chiral splitting}

Let $P_\tau:\mathfrak K\to\mathfrak K_\tau$ be the orthogonal projections and
define the two chiral response traces
\begin{equation}
 D_\tau(u)
 :=
 \tr_{\mathfrak K_\tau}\cM_u
 =
 2V_R\sum_{a=1}^3
 \|Q(K_a^\tau,u)\|_2^2.
 \label{R:rg:eq:Dtau}
\end{equation}
Then
\begin{equation}
       D_+(u)+D_-(u)=\langle Au,u\rangle.
 \label{R:rg:eq:Dsum}
\end{equation}

\begin{theorem}[Pure-shell chiral partition]
\label{R:rg:thm:pureshellsplit}
Let $u\in E_{k,\sigma}$ with $k\ge1$.  Then
\begin{equation}
 \boxed{
 D_\tau(u)
 =
 \frac{\kappa_k}{2}
 \left(
 1-\frac{2\tau\sigma}{k+2}
 \right)
 \|u\|_2^2.
 }
 \label{R:rg:eq:Dtauclosed}
\end{equation}
Equivalently,
\begin{align}
 D_\sigma(u)
 &=
 \frac{k}{2(k+2)}
 \kappa_k\|u\|_2^2,
 \label{R:rg:eq:same}\\
 D_{-\sigma}(u)
 &=
 \frac{k+4}{2(k+2)}
 \kappa_k\|u\|_2^2.
 \label{R:rg:eq:opposite}
\end{align}
Thus the opposite Killing chirality carries the larger partial response
trace, and the relative imbalance is
\begin{equation}
 \frac{
 D_{-\sigma}(u)-D_{\sigma}(u)
 }{
 D_{-\sigma}(u)+D_{\sigma}(u)
 }
 =
 \frac{2}{k+2}.
 \label{R:rg:eq:pureshellbias}
\end{equation}
\end{theorem}

\begin{proof}
Set
\[
 c_\tau(k,\sigma)
 :=
 \frac{k+2-2\tau\sigma}{k+2}.
\]
By \eqref{R:eq:KB},
\[
 Q(K_a^\tau,u)
 =
 \frac{c_\tau}{2}\mathcal L_{K_a^\tau}u.
\]
Using the factor Casimir \eqref{R:eq:partialCasimir},
\begin{align*}
 D_\tau(u)
 &=
 2V_R
 \frac{c_\tau^2}{4}
 \sum_{a=1}^3
 \|\mathcal L_{K_a^\tau}u\|_2^2\\
 &=
 \frac{2}{R^2}
 c_\tau^2
 j_\tau(j_\tau+1)
 \|u\|_2^2.
\end{align*}
If $\tau\sigma=1$, then
$c_\tau=k/(k+2)$ and $j_\tau=(k+2)/2$, giving
\eqref{R:rg:eq:same}.  If $\tau\sigma=-1$, then
$c_\tau=(k+4)/(k+2)$ and $j_\tau=k/2$, giving
\eqref{R:rg:eq:opposite}.  Formula \eqref{R:rg:eq:Dtauclosed} combines the two cases.
\end{proof}

\begin{remark}[High-frequency balance without connection flattening]
The relative trace bias \eqref{R:rg:eq:pureshellbias} tends to zero as
$k\to\infty$.  This does not contradict persistence of nonzero
operator-norm curvature in the shell-precession connection: response-trace
balance and connection flatness are different operator statements.
\end{remark}

\begin{theorem}[Global chiral trace identity]
\label{R:rg:thm:globalchiral}
Let $u$ be divergence-free with $\langle Au,u\rangle<\infty$.  Then
\begin{align}
 D_\tau(u)
 &=
 \sum_{k\ge1,\ \sigma=\pm}
 \frac{\kappa_k}{2}
 \left(
 1-\frac{2\tau\sigma}{k+2}
 \right)
 \|u_{k,\sigma}\|_2^2,
 \label{R:rg:eq:Dtauglobal}\\
 D_-(u)-D_+(u)
 &=
 \frac{2}{R}
 \langle A\curl^{-1}u,u\rangle.
 \label{R:rg:eq:chiraltraceidentity}
\end{align}
The Killing component does not contribute because $\kappa_0=0$.
\end{theorem}

\begin{proof}
Apply \cref{R:rg:thm:additivity,R:rg:thm:pureshellsplit}.  For the difference,
\[
 D_--D_+
 =
 2\sum_{k,\sigma}
 \sigma\frac{\kappa_k}{k+2}
 \|u_{k,\sigma}\|_2^2.
\]
Since
\[
 \frac{1}{\mu_{k,\sigma}}
 =
 \frac{\sigma R}{k+2},
\]
the right-hand side equals
\[
 \frac{2}{R}
 \sum_{k,\sigma}
 \frac{\kappa_k}{\mu_{k,\sigma}}
 \|u_{k,\sigma}\|_2^2
 =
 \frac{2}{R}
 \langle A\curl^{-1}u,u\rangle.
\]
\end{proof}

\begin{definition}[Chiral response index]
For $u$ with $\langle Au,u\rangle>0$, define
\begin{equation}
 \Xi(u)
 :=
 \frac{D_-(u)-D_+(u)}
      {D_-(u)+D_+(u)}.
 \label{R:rg:eq:Xi}
\end{equation}
\end{definition}

\begin{corollary}[Sharp universal chiral bound]
\label{R:rg:cor:Xibound}
For every non-Killing $u$ in the form domain of $A$,
\begin{equation}
 \boxed{
 |\Xi(u)|
 \le
 \frac{2}{3}.
 }
 \label{R:rg:eq:Xibound}
\end{equation}
The constant $2/3$ is sharp.  For $u\in E_{k,\sigma}$,
\begin{equation}
 \boxed{
 \Xi(u)=\frac{2\sigma}{k+2}.
 }
 \label{R:rg:eq:Xipureshell}
\end{equation}
\end{corollary}

\begin{proof}
By \eqref{R:rg:eq:Dtauglobal}, $\Xi$ is a convex weighted average of the numbers
$2\sigma/(k+2)$ with nonnegative weights
$\kappa_k\|u_{k,\sigma}\|_2^2$.  Since $k\ge1$ on the dissipative sector,
their absolute values are at most $2/3$.  Equality is attained on
$E_{1,+}$ or $E_{1,-}$.
\end{proof}

\subsection{The cross-chiral block}

Relative to
$\mathfrak K=\mathfrak K_+\oplus\mathfrak K_-$, write
\begin{equation}
 \cM_u
 =
 \begin{pmatrix}
 M^{++}_u & M^{+-}_u\\
 M^{-+}_u & M^{--}_u
 \end{pmatrix}.
 \label{R:rg:eq:blockM}
\end{equation}
The diagonal traces are $D_\pm(u)$.

\begin{proposition}[Universal cross-block bound]
\label{R:rg:prop:crossbound}
For every $u$,
\begin{equation}
 \boxed{
 \|M^{+-}_u\|_{\HS}
 \le
 \sqrt{D_+(u)D_-(u)}.
 }
 \label{R:rg:eq:crossbound}
\end{equation}
For a pure shell $u\in E_{k,\sigma}$,
\begin{equation}
 \|M^{+-}_u\|_{\HS}
 \le
 \frac{\sqrt{k(k+4)}}{2(k+2)}
 \kappa_k\|u\|_2^2.
 \label{R:rg:eq:crosspure}
\end{equation}
\end{proposition}

\begin{proof}
Split the response map as
\[
       \cB_u=(\cB_{u,+},\cB_{u,-})
\]
on the two Killing factors.  Then
\[
       M^{+-}_u=\cB_{u,+}^*\cB_{u,-}.
\]
Therefore
\[
 \|M^{+-}_u\|_{\HS}
 \le
 \|\cB_{u,+}\|_{\HS}
 \|\cB_{u,-}\|_{\HS}
 =
 \sqrt{D_+(u)D_-(u)}.
\]
Insert \cref{R:rg:thm:pureshellsplit} for \eqref{R:rg:eq:crosspure}.
\end{proof}

\begin{remark}
The cross block generally need not vanish for a fixed pure-shell vector.
Thus it is incorrect to infer from shell irreducibility that the full
$6\times6$ response tensor is block diagonal or scalar on each
$\mathfrak K_\tau$ for every $u$.  Only the two block traces are fixed
pointwise; full block isotropy emerges after an invariant ensemble average.
\end{remark}

\subsection{Uniform pure-shell ensemble}

Fix a shell $E=E_{k,\sigma}$ of real dimension $d_k$ and let $u$ be uniformly
distributed on its $L^2$ unit sphere.

\begin{theorem}[Ensemble-averaged response tensor]
\label{R:rg:thm:ensemble}
Let
\begin{equation}
 d_\tau^{(k,\sigma)}
 :=
 \frac{\kappa_k}{2}
 \left(
 1-\frac{2\tau\sigma}{k+2}
 \right).
 \label{R:rg:eq:dtaudef}
\end{equation}
Then
\begin{align}
 \mathbb E\,M^{\tau\tau}_u
 &=
 \frac{d_\tau^{(k,\sigma)}}{3}\Id_{\mathfrak K_\tau},
 \label{R:rg:eq:ensembleblock}\\
 \mathbb E\,M^{+-}_u
 &=
 0,
 \label{R:rg:eq:ensemblecross}\\
 \boxed{
 \mathbb E\,\cM_u
 =
 \begin{pmatrix}
 \frac{d_+^{(k,\sigma)}}{3}I_3 & 0\\
 0 & \frac{d_-^{(k,\sigma)}}{3}I_3
 \end{pmatrix}.
 }
 \label{R:rg:eq:ensembleM}
\end{align}
\end{theorem}

\begin{proof}
The uniform measure on the unit sphere of $E$ is invariant under the full
orthogonal group of $E$.  For fixed linear operators $S,T$ on $E$,
\[
 \mathbb E\,\langle Su,Tu\rangle
 =
 \frac{1}{d_k}\tr(S^*T).
 \label{R:rg:eq:sphereaverage}
\]
Within one simple factor, the trace pairing of the three representation
generators is an invariant bilinear form and therefore proportional to
$\delta_{ab}$.  Its total trace is already fixed by
\cref{R:rg:thm:pureshellsplit}, giving \eqref{R:rg:eq:ensembleblock}.

For opposite simple factors, the representation is a tensor product.
Each Lie generator is traceless on its irreducible spin factor, so
\[
 \tr\!\left[
 (\mathcal L_{K_a^+})^*
 \mathcal L_{K_b^-}
 \right]=0.
\]
Equation \eqref{R:rg:eq:sphereaverage} then gives
\eqref{R:rg:eq:ensemblecross}.
\end{proof}

\begin{remark}
The deterministic shell bias survives ensemble averaging through unequal
block traces, while directional anisotropy inside each
$\mathfrak K_\tau$ disappears in the mean.
\end{remark}

\subsection{Response polarization diagnostics}

The Casimir identity determines only the sum of the six eigenvalues.
To distinguish isotropic from concentrated response spectra, define
\begin{equation}
 p_i(u)
 :=
 \frac{\lambda_i(u)}{\tr\cM_u},
 \qquad
 \sum_{i=1}^6p_i=1,
 \label{R:rg:eq:pi}
\end{equation}
whenever $\tr\cM_u>0$.

Two useful invariants are
\begin{align}
 \mathcal P_\infty(u)
 &:=
 6p_1(u)
 =
 \frac{6\lambda_1(u)}{\tr\cM_u},
 \label{R:rg:eq:Pinfty}\\
 r_{\rm eff}(u)
 &:=
 \frac{1}{\sum_i p_i(u)^2}
 =
 \frac{(\tr\cM_u)^2}{\tr(\cM_u^2)}.
 \label{R:rg:eq:reff}
\end{align}

\begin{proposition}[Response-rank bounds]
\label{R:rg:prop:responsebounds}
For every non-Killing $u$,
\begin{equation}
 1\le\mathcal P_\infty(u)\le6,
 \qquad
 1\le r_{\rm eff}(u)\le6.
 \label{R:rg:eq:responsebounds}
\end{equation}
Both quantities are invariant under $SO(4)$ and under nonzero scalar
multiplication of $u$.
\end{proposition}

\begin{proof}
The statements are the standard inequalities for a probability vector
$(p_1,\ldots,p_6)$, together with
\cref{R:rg:thm:covariance} and $\cM_{cu}=c^2\cM_u$.
\end{proof}

\begin{remark}
A divergence of $\tr\cM_u$ says nothing by itself about the limiting values
of $\mathcal P_\infty$ or $r_{\rm eff}$.  Therefore the Casimir identity
forces at least one response to become large if dissipation becomes large,
but it does not force the normalized response spectrum to polarize.
\end{remark}

\subsection{Relation to helicity}

The chiral response index \eqref{R:rg:eq:Xi} is not the normalized fluid helicity.
Indeed,
\begin{equation}
 \mathscr H(u)
 =
 \langle u,\curl u\rangle
 =
 \sum_{k,\sigma}
 \sigma\frac{k+2}{R}
 \|u_{k,\sigma}\|_2^2,
 \label{R:rg:eq:helicity}
\end{equation}
whereas \eqref{R:rg:eq:chiraltraceidentity} uses the inverse-curl weight
$\kappa_k/\mu_{k,\sigma}$.

The two quantities emphasize opposite spectral ends: helicity weights high
curl frequency more strongly, while the response difference contains the
factor $1/(k+2)$ relative to the Stokes dissipation weight.

\begin{proposition}[Killing precession is transparent to weighted chirality]
\label{R:rg:prop:precessiontransparent}
Let $D(m)=2Q(m,\cdot)$ be the shellwise Killing-precession operator.  Then
\begin{equation}
 \langle D(m)v,A\curl^{-1}v\rangle=0
 \label{R:rg:eq:weightedtransparent}
\end{equation}
for every sufficiently regular $v$.
Equivalently, the quadratic form on the right side of
\eqref{R:rg:eq:chiraltraceidentity} receives no instantaneous contribution from
Killing precession.
\end{proposition}

\begin{proof}
The operator $D(m)$ is skew-adjoint and commutes with both $A$ and $\curl$
on the signed-shell decomposition.  Therefore
$A\curl^{-1}$ is self-adjoint and commutes with $D(m)$, so
\[
\begin{aligned}
 \langle Dv,A\curl^{-1}v\rangle
 &=
 -\langle v,D A\curl^{-1}v\rangle\\
 &=
 -\langle v,A\curl^{-1}Dv\rangle.
\end{aligned}
\]
which is the negative of the original real scalar.
\end{proof}

\begin{remark}
This transparency does not imply that the nonlinear high--high term
$Q(v,v)$ preserves shellwise $+$/$-$ energy equality.  In particular, zero
initial helicity does not force pairwise equality
$\|v_{k,+}(t)\|_2=\|v_{k,-}(t)\|_2$ for later times.
\end{remark}


\section{Stokes response moments and finite-band tomography}

\subsection{The Casimir response hierarchy}

Because $A$ is invertible on $\cK^\perp$, negative fractional powers are
well defined there.

\begin{definition}[Response hierarchy]
\label{R:hi:def:hierarchy}
Let $s\ge-1$ and suppose
$v\in\cK^\perp$ belongs to the domain of $A^{(s+1)/2}$.
For $X,Y\in\cK$, define
\begin{equation}
 \boxed{
 \cM_v^{[s]}(X,Y)
 :=
 2V_R
 \left\langle
 A^{s/2}Q(X,v),
 A^{s/2}Q(Y,v)
 \right\rangle.
 }
 \label{R:hi:eq:Ms}
\end{equation}
For $s=-1$, $A^{-1/2}$ is taken on $\cK^\perp$.
\end{definition}

Each $\cM_v^{[s]}$ is a positive semidefinite symmetric operator on the
six-dimensional Killing algebra.

\begin{theorem}[Exact trace ladder]
\label{R:hi:thm:traceladder}
For every $s\ge-1$,
\begin{equation}
 \boxed{
 \tr_{\cK}\cM_v^{[s]}
 =
 \langle A^{s+1}v,v\rangle
 =
 \left\|A^{(s+1)/2}v\right\|_2^2.
 }
 \label{R:hi:eq:traceladder}
\end{equation}
In particular,
\begin{align}
 \tr\cM_v^{[-1]}
 &=
 \|v\|_2^2,
 \label{R:hi:eq:energyresponse}\\
 \tr\cM_v^{[0]}
 &=
 \langle Av,v\rangle
 =
 2\|\Def v\|_2^2.
 \label{R:hi:eq:dissresponse}
\end{align}
\end{theorem}

\begin{proof}
Since every $T_a^\tau$ commutes with $A$ and is skew-adjoint,
\begin{align*}
 \tr\cM_v^{[s]}
 &=
 2V_R
 \sum_{\tau,a}
 \left\langle
 T_a^\tau v,A^sT_a^\tau v
 \right\rangle\\
 &=
 -2V_R
 \sum_{\tau,a}
 \left\langle
 v,A^s(T_a^\tau)^2v
 \right\rangle.
\end{align*}
Use \eqref{R:eq:CasimirCore}.
\end{proof}

\begin{theorem}[Spectral additivity]
\label{R:hi:thm:additivity}
If
\[
       v=\sum_{k\ge1,\ \sigma=\pm}v_{k,\sigma},
\]
then
\begin{equation}
 \boxed{
 \cM_v^{[s]}
 =
 \sum_{k,\sigma}
 \cM_{v_{k,\sigma}}^{[s]}.
 }
 \label{R:hi:eq:additivity}
\end{equation}
Whenever the trace in \eqref{R:hi:eq:traceladder} is finite, the series converges
in operator norm on $\cK$.
\end{theorem}

\begin{proof}
Each $Q(X,v_{k,\sigma})$ remains in $E_{k,\sigma}$, and distinct signed
shells are orthogonal.  Hence all cross-shell Gram terms vanish.  The
summands are positive semidefinite and their traces sum to
\eqref{R:hi:eq:traceladder}; finite-dimensional monotone convergence then gives
operator-norm convergence.
\end{proof}

\subsection{The Stokes moment sequence}

Define the total energy at Stokes level $k$ by
\begin{equation}
 E_k
 :=
 \|v_{k,+}\|_2^2+\|v_{k,-}\|_2^2.
 \label{R:hi:eq:Ek}
\end{equation}

\begin{definition}[Casimir response moments]
For integers $n\ge0$, set
\begin{equation}
 \boxed{
 m_n(v)
 :=
 \tr\cM_v^{[n-1]}
 =
 \langle A^n v,v\rangle.
 }
 \label{R:hi:eq:mn}
\end{equation}
\end{definition}

By the signed-shell decomposition,
\begin{equation}
 \boxed{
 m_n
 =
 \sum_{k\ge1}\kappa_k^nE_k.
 }
 \label{R:hi:eq:moments}
\end{equation}
Thus $(m_n)$ is the moment sequence of the positive discrete measure
\begin{equation}
 \mu_v
 :=
 \sum_{k\ge1}E_k\,\delta_{\kappa_k}.
 \label{R:hi:eq:measure}
\end{equation}

\begin{theorem}[Hankel positivity and exact spectral rank]
\label{R:hi:thm:Hankel}
For $N\ge0$, define the Hankel matrix
\begin{equation}
 H_N(v)
 :=
 [m_{i+j}(v)]_{i,j=0}^N.
 \label{R:hi:eq:Hankel}
\end{equation}
Then $H_N(v)$ is positive semidefinite.

If $v$ has exactly $r$ active Stokes levels
$k_1,\ldots,k_r$, then
\begin{equation}
 \boxed{
 \operatorname{rank}H_N(v)
 =
 \min\{N+1,r\}.
 }
 \label{R:hi:eq:Hankelrank}
\end{equation}
Consequently,
\begin{equation}
 \det H_{r-1}>0,
 \qquad
 \det H_r=0.
 \label{R:hi:eq:detcriterion}
\end{equation}
\end{theorem}

\begin{proof}
For $c=(c_0,\ldots,c_N)$ and
$p(\lambda)=\sum_{j=0}^Nc_j\lambda^j$,
\begin{align*}
 c^\mathsf T H_Nc
 &=
 \sum_{i,j}c_ic_jm_{i+j}\\
 &=
 \sum_kE_k|p(\kappa_k)|^2
 \ge0.
\end{align*}
If the support consists of $r$ distinct nodes, write
\[
 H_N
 =
 V_N
 \operatorname{diag}(E_{k_1},\ldots,E_{k_r})
 V_N^\mathsf T,
\]
where $(V_N)_{ij}=\kappa_{k_j}^i$.
The Vandermonde matrix has rank $\min\{N+1,r\}$.
\end{proof}

\begin{remark}
Hankel rank is used broadly in finite moment problems and realization theory
\cite{Fasino1995}.  Here the point is not the linear-algebra fact itself, but
that the required moments are exact traces of Killing-response tensors.
\end{remark}

\subsection{Log-convexity and an exact spectral-spread identity}

\begin{theorem}[Log-convex response moments]
\label{R:hi:thm:logconvex}
For every $n\ge1$ for which the moments are finite,
\begin{equation}
 \boxed{
 m_n^2
 \le
 m_{n-1}m_{n+1}.
 }
 \label{R:hi:eq:logconvex}
\end{equation}
If $v\ne0$, equality holds for some $n\ge1$ if and only if the non-Killing
field is supported on one Stokes level
\[
       E_{k,+}\oplus E_{k,-}
\]
for a single $k$.
\end{theorem}

\begin{proof}
Apply Cauchy--Schwarz to
$A^{(n-1)/2}v$ and $A^{(n+1)/2}v$.
Equality requires these two vectors to be linearly dependent.  Since $A$ is
self-adjoint with distinct eigenvalues $\kappa_k$ for distinct levels, this
occurs exactly on one Stokes eigenspace.
\end{proof}

For the first inequality one can write the defect exactly.

\begin{theorem}[Spectral-variance identity]
\label{R:hi:thm:variance}
For every $v\in D(A)$,
\begin{equation}
 \boxed{
 m_0m_2-m_1^2
 =
 \frac12
 \sum_{k,\ell\ge1}
 E_kE_\ell
 (\kappa_k-\kappa_\ell)^2.
 }
 \label{R:hi:eq:variance}
\end{equation}
Equivalently,
\begin{equation}
 \frac{m_2}{m_0}
 -
 \left(\frac{m_1}{m_0}\right)^2
 =
 \operatorname{Var}_{p}(\kappa),
 \qquad
 p_k=\frac{E_k}{m_0}.
 \label{R:hi:eq:normalizedvariance}
\end{equation}
Thus three response traces detect exact Stokes monochromaticity.
\end{theorem}

\begin{proof}
Expand the right side:
\[
 \frac12\sum_{k,\ell}E_kE_\ell
 (\kappa_k^2+\kappa_\ell^2-2\kappa_k\kappa_\ell)
 =
 m_0m_2-m_1^2.
\]
\end{proof}

\begin{corollary}[Monotone effective Stokes frequency]
\label{R:hi:cor:effectivefrequency}
Assume all moments involved are finite and define
\begin{equation}
 \Omega_n(v)
 :=
 \frac{m_{n+1}(v)}{m_n(v)}.
 \label{R:hi:eq:Omega}
\end{equation}
Then $\Omega_n$ is nondecreasing in $n$.  For finite spectral support,
\begin{equation}
 \lim_{n\to\infty}\Omega_n
 =
 \max\{\kappa_k:E_k>0\}.
 \label{R:hi:eq:Omegalimit}
\end{equation}
\end{corollary}

\begin{proof}
Monotonicity is \eqref{R:hi:eq:logconvex}.  For finite support, divide numerator
and denominator by the largest active eigenvalue to the power $n$.
\end{proof}

\subsection{Chiral moment sequence}

Let
\begin{equation}
 D_\tau^{[s]}(v)
 :=
 \tr_{\cK_\tau}\cM_v^{[s]}.
 \label{R:hi:eq:Dtaus}
\end{equation}
The signed-shell Killing--Beltrami algebra gives the pure-shell formula
\begin{equation}
 D_\tau^{[s]}(v_{k,\sigma})
 =
 \frac{\kappa_k^{s+1}}{2}
 \left(
 1-\frac{2\tau\sigma}{k+2}
 \right)
 \|v_{k,\sigma}\|_2^2.
 \label{R:hi:eq:pureshellD}
\end{equation}

\begin{theorem}[Higher chiral trace identity]
\label{R:hi:thm:chiralhierarchy}
For every $s\ge-1$,
\begin{align}
 D_+^{[s]}(v)+D_-^{[s]}(v)
 &=
 \langle A^{s+1}v,v\rangle,
 \label{R:hi:eq:chiralsum}\\
 \boxed{
 D_-^{[s]}(v)-D_+^{[s]}(v)
 =
 \frac{2}{R}
 \langle A^{s+1}\curl^{-1}v,v\rangle.
 }
 \label{R:hi:eq:chiraldiff}
\end{align}
\end{theorem}

\begin{proof}
Sum \eqref{R:hi:eq:pureshellD} over signed shells and use
\[
       \mu_{k,\sigma}^{-1}
       =
       \sigma\frac{R}{k+2}.
\]
\end{proof}

\begin{definition}[Signed response moments]
For $n\ge0$, define
\begin{equation}
\begin{aligned}
 h_n(v)
 &:=
 \frac{R}{2}
 \left(
 D_-^{[n-1]}(v)-D_+^{[n-1]}(v)
 \right)\\
 &=
 \langle A^n\curl^{-1}v,v\rangle.
\end{aligned}
\label{R:hi:eq:hn}
\end{equation}
\end{definition}

Writing
\[
 E_{k,\pm}
 :=
 \|v_{k,\pm}\|_2^2,
\]
one obtains
\begin{equation}
 \boxed{
 h_n
 =
 \sum_{k\ge1}
 \kappa_k^n
 \frac{R}{k+2}
 \bigl(E_{k,+}-E_{k,-}\bigr).
 }
 \label{R:hi:eq:hmoments}
\end{equation}

\subsection{Finite-band spectral tomography}

We now show that the two moment sequences recover the complete signed-shell
energy spectrum.

\begin{theorem}[Exact finite-band tomography]
\label{R:hi:thm:tomography}
Suppose $v$ has nonzero energy on exactly $r$ distinct Stokes levels
\[
       k_1,\ldots,k_r.
\]
Then the finite data
\begin{equation}
 m_0,m_1,\ldots,m_{2r},
 \qquad
 h_0,h_1,\ldots,h_{r-1}
 \label{R:hi:eq:tomodata}
\end{equation}
determine uniquely
\[
       \{k_j,E_{k_j,+},E_{k_j,-}\}_{j=1}^r.
\]
\end{theorem}

\begin{proof}
By \cref{R:hi:thm:Hankel}, $r$ is the unique integer for which
\[
       \det H_{r-1}>0,
       \qquad
       \det H_r=0.
\]

Let
\[
 p(z)
 =
 \prod_{j=1}^r(z-\kappa_{k_j})
 =
 z^r+c_{r-1}z^{r-1}+\cdots+c_0.
\]
Since $p(\kappa_{k_j})=0$,
\begin{equation}
 \sum_{\ell=0}^r
 c_\ell m_{n+\ell}
 =
 0,
 \qquad
 n=0,\ldots,r-1,
 \qquad
 c_r=1.
 \label{R:hi:eq:recurrence}
\end{equation}
The coefficient vector $(c_0,\ldots,c_{r-1})$ therefore solves
\begin{equation}
 H_{r-1}
 \begin{pmatrix}
 c_0\\ \vdots\\ c_{r-1}
 \end{pmatrix}
 =
 -
 \begin{pmatrix}
 m_r\\ \vdots\\ m_{2r-1}
 \end{pmatrix}.
 \label{R:hi:eq:annihilator}
\end{equation}
The matrix is positive definite, so the solution is unique.  The roots of
$p$ are exactly the active Stokes eigenvalues.  Since
$\kappa_k=k(k+4)/R^2$ is strictly increasing in $k$, the levels $k_j$ are
recovered uniquely.

Now solve the Vandermonde system
\begin{equation}
 m_n
 =
 \sum_{j=1}^r
 \kappa_{k_j}^n E_{k_j},
 \qquad
 n=0,\ldots,r-1,
 \label{R:hi:eq:totalVandermonde}
\end{equation}
for
\[
       E_{k_j}=E_{k_j,+}+E_{k_j,-}.
\]
Likewise solve
\begin{equation}
 h_n
 =
 \sum_{j=1}^r
 \kappa_{k_j}^n\delta_j,
 \qquad
 \delta_j
 :=
 \frac{R}{k_j+2}
 (E_{k_j,+}-E_{k_j,-}),
 \label{R:hi:eq:signedVandermonde}
\end{equation}
for $n=0,\ldots,r-1$.  Finally,
\begin{equation}
 \boxed{
 E_{k_j,\pm}
 =
 \frac12
 \left[
 E_{k_j}
 \pm
 \frac{k_j+2}{R}\delta_j
 \right].
 }
 \label{R:hi:eq:recoverenergies}
\end{equation}
\end{proof}

\begin{remark}
The reconstruction is an exact finite moment problem and can be implemented
by standard Prony/Hankel techniques.  Numerical conditioning is a separate
issue; Hankel moment matrices are well known to become ill-conditioned at
high order \cite{Fasino1995}.
\end{remark}

\subsection{Complete monotonicity under Stokes relaxation}

Consider the unforced linear Stokes equation on $\cK^\perp$,
\begin{equation}
       \partial_tv+\nu Av=0.
       \label{R:hi:eq:Stokes}
\end{equation}

\begin{theorem}[Loewner-complete monotonicity]
\label{R:hi:thm:matrixCM}
For every $s\ge-1$,
\begin{equation}
 \boxed{
 \frac{\dd}{\dd t}
 \cM_{v(t)}^{[s]}
 =
 -2\nu
 \cM_{v(t)}^{[s+1]}.
 }
 \label{R:hi:eq:matrixladder}
\end{equation}
Consequently, for every integer $q\ge0$,
\begin{equation}
 \boxed{
 (-1)^q
 \frac{\dd^q}{\dd t^q}
 \cM_{v(t)}^{[s]}
 =
 (2\nu)^q
 \cM_{v(t)}^{[s+q]}
 \succeq0.
 }
 \label{R:hi:eq:matrixCM}
\end{equation}
Thus $t\mapsto\cM_{v(t)}^{[s]}$ is completely monotone in the Loewner
order.
\end{theorem}

\begin{proof}
Since every Killing response commutes with $A$,
\[
 \frac{\dd}{\dd t}
 A^{s/2}Q(X,v)
 =
 -\nu
 A^{s/2+1}Q(X,v).
\]
Differentiate the two factors in \eqref{R:hi:eq:Ms}.  The two equal terms combine
to give \eqref{R:hi:eq:matrixladder}; iteration gives \eqref{R:hi:eq:matrixCM}.
\end{proof}

\begin{corollary}[Scalar moment derivatives]
\label{R:hi:cor:mderivatives}
Under \eqref{R:hi:eq:Stokes},
\begin{equation}
 \frac{\dd}{\dd t}m_n(t)
 =
 -2\nu m_{n+1}(t),
 \label{R:hi:eq:mderivative}
\end{equation}
and
\begin{equation}
 (-1)^q
 \frac{\dd^q}{\dd t^q}
 m_n(t)
 =
 (2\nu)^q
 m_{n+q}(t).
 \label{R:hi:eq:scalarCM}
\end{equation}
\end{corollary}

\subsection{Two-signal Stokes spectroscopy}

Let
\begin{align}
 \mathcal E(t)
 &:=
 \|v(t)\|_2^2,
 \label{R:hi:eq:Etime}\\
 \mathcal C(t)
 &:=
 \langle\curl^{-1}v(t),v(t)\rangle.
 \label{R:hi:eq:Ctime}
\end{align}

\begin{theorem}[Relaxation tomography]
\label{R:hi:thm:relaxationtomography}
For the Stokes evolution \eqref{R:hi:eq:Stokes},
\begin{align}
 \mathcal E(t)
 &=
 \sum_{k\ge1}
 E_k(0)e^{-2\nu\kappa_kt},
 \label{R:hi:eq:Eexp}\\
 \mathcal C(t)
 &=
 \sum_{k\ge1}
 \frac{R}{k+2}
 \bigl(E_{k,+}(0)-E_{k,-}(0)\bigr)
 e^{-2\nu\kappa_kt}.
 \label{R:hi:eq:Cexp}
\end{align}
Moreover,
\begin{align}
 \boxed{
 m_n(0)
 =
 \frac{(-1)^n}{(2\nu)^n}
 \mathcal E^{(n)}(0),
 }
 \label{R:hi:eq:mfromE}\\
 \boxed{
 h_n(0)
 =
 \frac{(-1)^n}{(2\nu)^n}
 \mathcal C^{(n)}(0).
 }
 \label{R:hi:eq:hfromC}
\end{align}
Hence, for a finite-band state with $r$ active Stokes levels, the derivatives
\[
 \mathcal E^{(0)}(0),\ldots,\mathcal E^{(2r)}(0),
 \qquad
 \mathcal C^{(0)}(0),\ldots,\mathcal C^{(r-1)}(0)
\]
determine the complete signed-shell energy distribution.
\end{theorem}

\begin{proof}
The Stokes semigroup multiplies each shell by
$e^{-\nu\kappa_kt}$.  Squaring gives
\eqref{R:hi:eq:Eexp}--\eqref{R:hi:eq:Cexp}.  Differentiate at $t=0$ and use
\cref{R:hi:thm:tomography}.
\end{proof}

\begin{remark}[Relation to Hankel methods in fluid mechanics]
Time-delay Hankel matrices are widely used for data-driven fluid
decompositions and are closely related to space-time POD
\cite{FrameTowne2023}.  The two scalar curves above solve a different
inverse problem: in the exact round-sphere Stokes dynamics their exponential
nodes are fixed geometric eigenvalues, and their amplitudes are the signed
shell energies.
\end{remark}

\subsection{The nonlinear tensor balance}

Return to the high-mode Navier--Stokes equation on an affine Killing leaf,
\begin{equation}
 \partial_tv
 +
 \nu Av
 +
 Q(v,v)
 +
 D(m)v
 =
 F_\perp,
 \label{R:hi:eq:highNS}
\end{equation}
where $D(m)$ is skew-adjoint and commutes with $A$.

For $s\ge-1$ and high-mode fields $v,w$, define the symmetric cross-response
tensor
\begin{align}
 \cC_s[v;w](X,Y)
 &:=
 2V_R
 \Bigl[
 \langle
 A^{s/2}Q(X,w),
 A^{s/2}Q(Y,v)
 \rangle
 \notag\\
 &\qquad\qquad+
 \langle
 A^{s/2}Q(X,v),
 A^{s/2}Q(Y,w)
 \rangle
 \Bigr].
 \label{R:hi:eq:Cs}
\end{align}

\begin{lemma}[Trace polarization]
\label{R:hi:lem:tracepolarization}
For all sufficiently regular $v,w$,
\begin{equation}
 \boxed{
 \tr\cC_s[v;w]
 =
 2\langle A^{s+1}v,w\rangle.
 }
 \label{R:hi:eq:tracepolarization}
\end{equation}
\end{lemma}

\begin{proof}
Polarize the proof of \cref{R:hi:thm:traceladder}.
\end{proof}

\begin{theorem}[Exact nonlinear response balance]
\label{R:hi:thm:nonlinearbalance}
Along \eqref{R:hi:eq:highNS},
\begin{equation}
\begin{aligned}
 \frac{\dd}{\dd t}\cM_v^{[s]}
 &=
 -2\nu\cM_v^{[s+1]}
 -
 \cC_s[v;Q(v,v)]\\
 &\quad-
 \cC_s[v;D(m)v]
 +
 \cC_s[v;F_\perp].
\end{aligned}
\label{R:hi:eq:tensorbalance}
\end{equation}
Taking the trace gives, with $n=s+1\ge0$,
\begin{equation}
\begin{aligned}
 \frac{\dd}{\dd t}\langle A^n v,v\rangle
 &=
 -2\nu\langle A^{n+1}v,v\rangle\\
 &\quad-
 2\langle A^n v,Q(v,v)\rangle
 +
 2\langle A^n v,F_\perp\rangle.
\end{aligned}
\label{R:hi:eq:momentbalance}
\end{equation}
The Killing-precession term has zero trace at every level.
\end{theorem}

\begin{proof}
Differentiate \eqref{R:hi:eq:Ms} and substitute \eqref{R:hi:eq:highNS}.  Since
$Q(X,\cdot)$ commutes with $A$, the viscous contribution is exactly
$-2\nu\cM^{[s+1]}$.  Taking the trace and using
\cref{R:hi:lem:tracepolarization} gives \eqref{R:hi:eq:momentbalance}.  Finally,
\[
 \langle A^n v,D(m)v\rangle=0
\]
because $D(m)$ is skew-adjoint and commutes with $A^n$.
\end{proof}

\begin{corollary}[Trace-free nonlinear redistribution at the energy level]
\label{R:hi:cor:tracefree}
At $s=-1$,
\begin{equation}
 \tr\cC_{-1}[v;Q(v,v)]
 =
 2\langle v,Q(v,v)\rangle
 =
 0.
 \label{R:hi:eq:tracefreeQ}
\end{equation}
Also,
\begin{equation}
 \tr\cC_{-1}[v;D(m)v]=0.
 \label{R:hi:eq:tracefreeD}
\end{equation}
Hence, in the unforced equation, the energy-response tensor obeys
\begin{equation}
 \frac{\dd}{\dd t}\cM_v^{[-1]}
 =
 -2\nu\cM_v^{[0]}
 +
 \mathcal R_{\rm tf}(v,m),
 \qquad
 \tr\mathcal R_{\rm tf}=0,
 \label{R:hi:eq:energymatrixbalance}
\end{equation}
where the trace-free term redistributes the six response channels without
changing their total high-mode energy.
\end{corollary}

\begin{proof}
The first identity is the standard skew-energy property of incompressible
convection.  The second follows from skew-adjointness of $D(m)$.
\end{proof}

\begin{remark}
Trace-free does not mean norm-free or dynamically negligible.  The nonlinear
and Killing terms can substantially change the eigenvectors and eigenvalue
distribution of $\cM^{[-1]}$ while preserving its trace instantaneously.
\end{remark}

\subsection{A response-moment notion of spectral complexity}

The Hankel hierarchy gives an exact algebraic notion of the number of active
Stokes levels.

\begin{definition}
For a finite-band nonzero field, define its response spectral rank by
\begin{equation}
 r_{\rm spec}(v)
 :=
 \operatorname{rank}
 [m_{i+j}]_{i,j\ge0}.
 \label{R:hi:eq:rspec}
\end{equation}
\end{definition}

\begin{corollary}
For finite-band $v$,
\begin{equation}
 \boxed{
 r_{\rm spec}(v)
 =
 \#\{k\ge1:E_k>0\}.
 }
 \label{R:hi:eq:rspeclevels}
\end{equation}
This quantity is invariant under $SO(4)$ and under nonzero rescaling of
$v$.
\end{corollary}

The definition counts distinct Stokes eigenvalues, not the number of signed
shells or individual modes.  The signed moments $h_n$ provide the additional
chirality information required to split each active level.


\section{Curl-response Fourier tomography}

\subsection{Curl-response moments}

\begin{definition}[Curl-response tensor]
For an integer $n\ge0$ and sufficiently regular
$v\in\cK^\perp$, define the symmetric bilinear form on $\cK$
\begin{equation}
 \boxed{
 \cC_n(v)(X,Y)
 :=
 2V_R
 \left\langle
 Q(X,v),
 \curl^n Q(Y,v)
 \right\rangle.
 }
 \label{R:cf:eq:Cn}
\end{equation}
\end{definition}

Symmetry follows because $\curl^n$ is self-adjoint.

For a single signed shell define its zeroth response weight
\begin{equation}
 \cN_{k,\sigma}(v)
 :=
 \cC_0(v_{k,\sigma}).
 \label{R:cf:eq:Nks}
\end{equation}
It is a positive semidefinite operator on the six-dimensional Killing
algebra.

\begin{theorem}[Matrix-valued curl moment representation]
\label{R:cf:thm:momentrepresentation}
Let
\[
       v=\sum_{k\ge1,\sigma=\pm}v_{k,\sigma}
\]
and suppose the moment in question is finite.  Then
\begin{equation}
 \boxed{
 \cC_n(v)
 =
 \sum_{k\ge1,\sigma=\pm}
 \mu_{k,\sigma}^n
 \cN_{k,\sigma}(v).
 }
 \label{R:cf:eq:matrixmoments}
\end{equation}
Moreover
\begin{equation}
 \boxed{
 \tr\cN_{k,\sigma}(v)
 =
 \kappa_k\|v_{k,\sigma}\|_2^2
 }
 \label{R:cf:eq:Ntrace}
\end{equation}
and therefore
\begin{equation}
 \boxed{
 c_n(v)
 :=
 \tr\cC_n(v)
 =
 \langle A\curl^n v,v\rangle
 =
 \sum_{k,\sigma}
 \kappa_k\mu_{k,\sigma}^n
 \|v_{k,\sigma}\|_2^2.
 }
 \label{R:cf:eq:scalarcurlmoments}
\end{equation}
\end{theorem}

\begin{proof}
For every Killing field $X$,
\[
       Q(X,v_{k,\sigma})\in E_{k,\sigma}.
\]
Distinct signed shells are orthogonal.  Hence the Gram tensor splits over
shells, and applying $\curl^n$ to the shell response multiplies it by
$\mu_{k,\sigma}^n$, proving \eqref{R:cf:eq:matrixmoments}.

Taking $n=0$ and using \eqref{R:eq:CasimirCore},
\[
 \tr\cN_{k,\sigma}
 =
 2V_R\sum_{\tau,a}
 \|T_a^\tau v_{k,\sigma}\|_2^2
 =
 \langle Av_{k,\sigma},v_{k,\sigma}\rangle,
\]
which is \eqref{R:cf:eq:Ntrace}.

Finally, the same Casimir argument, using commutation with curl, gives
\[
 2V_R\sum_{\tau,a}
 \langle T_a^\tau v,\curl^nT_a^\tau v\rangle
 =
 \langle A\curl^n v,v\rangle.
\]
\end{proof}

\begin{definition}[Response spectral measure]
If $\langle Av,v\rangle<\infty$, define the positive
$\operatorname{Sym}^+(\cK)$-valued measure
\begin{equation}
 \boxed{
 \Sigma_v
 :=
 \sum_{k\ge1,\sigma=\pm}
 \cN_{k,\sigma}(v)\,
 \delta_{\mu_{k,\sigma}}.
 }
 \label{R:cf:eq:Sigmav}
\end{equation}
\end{definition}

Its total trace mass is finite:
\[
        \tr\Sigma_v(\mathbb R)
        =
        \langle Av,v\rangle.
\]
Equation \eqref{R:cf:eq:matrixmoments} is precisely the matrix moment formula
\[
       \cC_n
       =
       \int_{\mathbb R}\lambda^n\,d\Sigma_v(\lambda).
\]

\subsection{Block Hankel positivity and response-channel rank}

For $N\ge0$ form the $6(N+1)\times6(N+1)$ block Hankel matrix
\begin{equation}
 \mathbb H_N(v)
 :=
 [\cC_{i+j}(v)]_{i,j=0}^N.
 \label{R:cf:eq:blockHankel}
\end{equation}

\begin{theorem}[Block Hankel positivity]
\label{R:cf:thm:blockpositive}
Whenever the moments through order $2N$ are finite,
\begin{equation}
 \boxed{
 \mathbb H_N(v)\succeq0.
 }
 \label{R:cf:eq:blockpositive}
\end{equation}
\end{theorem}

\begin{proof}
For $X_0,\ldots,X_N\in\cK$,
\begin{align*}
 \sum_{i,j=0}^N
 \langle X_i,\cC_{i+j}X_j\rangle
 &=
 \sum_{k,\sigma}
 \left\langle
 \sum_i\mu_{k,\sigma}^iX_i,
 \cN_{k,\sigma}
 \sum_j\mu_{k,\sigma}^jX_j
 \right\rangle\\
 &\ge0.
\end{align*}
\end{proof}

The scalar Hankel matrix
\[
       H_N^{\rm tr}
       =
       [c_{i+j}]_{i,j=0}^N
\]
is therefore positive as well.

\begin{theorem}[Finite-support rank laws]
\label{R:cf:thm:ranklaws}
Suppose exactly $q$ signed shells
$\alpha=1,\ldots,q$ are active, with distinct curl eigenvalues
$\mu_\alpha$ and nonzero response weights $\cN_\alpha$.  Then
\begin{equation}
 \boxed{
 \operatorname{rank}H_N^{\rm tr}
 =
 \min\{N+1,q\}.
 }
 \label{R:cf:eq:scalarrank}
\end{equation}
For the block Hankel matrix, if $N\ge q-1$,
\begin{equation}
 \boxed{
 \operatorname{rank}\mathbb H_N
 =
 \sum_{\alpha=1}^q
 \operatorname{rank}\cN_\alpha.
 }
 \label{R:cf:eq:blockrank}
\end{equation}
\end{theorem}

\begin{proof}
The scalar weights
\[
       w_\alpha
       =
       \tr\cN_\alpha
       =
       \kappa_\alpha\|v_\alpha\|_2^2
\]
are strictly positive.  Thus
\[
 H_N^{\rm tr}
 =
 V_N
 \operatorname{diag}(w_1,\ldots,w_q)
 V_N^\mathsf T,
\]
where $(V_N)_{i\alpha}=\mu_\alpha^i$.  This gives
\eqref{R:cf:eq:scalarrank}.

Similarly,
\begin{equation}
 \mathbb H_N
 =
 \mathbb V_N
 \operatorname{diag}(\cN_1,\ldots,\cN_q)
 \mathbb V_N^*,
 \label{R:cf:eq:blockfactor}
\end{equation}
where $\mathbb V_N$ is the block Vandermonde operator with
$\alpha$th block column
\[
       (1,\mu_\alpha,\ldots,\mu_\alpha^N)^\mathsf T\otimes I_{\cK}.
\]
For $N\ge q-1$ the scalar Vandermonde columns are independent, so
$\mathbb V_N$ is injective on the direct sum of the $q$ copies of $\cK$.
Hence multiplication by $\mathbb V_N$ and its adjoint preserves the rank of
the block diagonal weight.
\end{proof}

\begin{remark}
The two ranks measure different complexity.  The scalar rank counts active
signed curl shells.  The block rank counts the sum of the independent
Killing-response channel ranks carried by those shells.
\end{remark}

\begin{corollary}[Block-Hankel determinant]
\label{R:cf:cor:blockdet}
Assume $N=q-1$ and every active $\cN_\alpha$ is positive definite on
$\cK$.  Then
\begin{equation}
 \boxed{
 \det\mathbb H_{q-1}
 =
 \left[
 \prod_{\alpha<\beta}
 (\mu_\beta-\mu_\alpha)
 \right]^{12}
 \prod_{\alpha=1}^q
 \det\cN_\alpha.
 }
 \label{R:cf:eq:blockdet}
\end{equation}
\end{corollary}

\begin{proof}
In \eqref{R:cf:eq:blockfactor}, $\mathbb V_{q-1}=V\otimes I_6$ with square
Vandermonde $V$.  Therefore
\[
\begin{aligned}
 \det\mathbb H_{q-1}
 &=
 \det(V\otimes I_6)^2
 \prod_\alpha\det\cN_\alpha\\
 &=
 \det(V)^{12}
 \prod_\alpha\det\cN_\alpha.
\end{aligned}
\]
\end{proof}

\subsection{Finite-band matrix recurrence}

\begin{theorem}[Annihilating recurrence]
\label{R:cf:thm:recurrence}
Under the finite-support assumptions of \cref{R:cf:thm:ranklaws}, let
\begin{equation}
 p(z)
 =
 \prod_{\alpha=1}^q(z-\mu_\alpha)
 =
 z^q+a_{q-1}z^{q-1}+\cdots+a_0.
 \label{R:cf:eq:annpoly}
\end{equation}
Then the full matrix sequence obeys
\begin{equation}
 \boxed{
 \cC_{n+q}
 +
 a_{q-1}\cC_{n+q-1}
 +\cdots+
 a_0\cC_n
 =
 0,
 \qquad
 n\ge0.
 }
 \label{R:cf:eq:matrixrecurrence}
\end{equation}
The same polynomial is the minimal scalar recurrence for
$c_n=\tr\cC_n$.
\end{theorem}

\begin{proof}
Insert \eqref{R:cf:eq:matrixmoments}.  The coefficient of each
$\cN_\alpha$ is
$\mu_\alpha^np(\mu_\alpha)=0$.
Minimality for the scalar sequence follows because all scalar weights are
strictly positive and the nodes are distinct.
\end{proof}

Thus the scalar trace moments identify the signed curl eigenvalues, while
the first $q$ matrix moments determine the full response weights through an
operator-valued Vandermonde system.

\subsection{The curl-response characteristic function}

The spectrum \eqref{G:eq:curlspec} lies in $(1/R)\mathbb Z$.  Hence
\begin{equation}
       e^{i2\pi R\curl}=\Id
 \label{R:cf:eq:periodiccurl}
\end{equation}
on divergence-free vector fields.

\begin{definition}[Operator-valued response characteristic function]
For $v$ with finite Stokes energy define
\begin{equation}
\begin{aligned}
 \Phi_v(\theta)(X,Y)
 &:=
 2V_R
 \left\langle
 Q(X,v),
 e^{i\theta\curl}Q(Y,v)
 \right\rangle,\\
 &\hspace{2em}\theta\in\mathbb R.
\end{aligned}
\label{R:cf:eq:Phi}
\end{equation}
\end{definition}

The function is $2\pi R$-periodic and satisfies
\[
       \Phi_v(-\theta)=\Phi_v(\theta)^*.
\]

\begin{theorem}[Exact operator-valued Fourier representation]
\label{R:cf:thm:Fourier}
For every finite-dissipation $v$,
\begin{equation}
 \boxed{
 \Phi_v(\theta)
 =
 \sum_{k\ge1,\sigma=\pm}
 e^{i\sigma(k+2)\theta/R}
 \cN_{k,\sigma}(v),
 }
 \label{R:cf:eq:Fourier}
\end{equation}
where the series converges absolutely in trace norm, hence uniformly in
operator norm.

Each shell-response tensor is recovered by Fourier inversion:
\begin{equation}
 \boxed{
 \cN_{k,\sigma}(v)
 =
 \frac{1}{2\pi R}
 \int_0^{2\pi R}
 \Phi_v(\theta)
 e^{-i\sigma(k+2)\theta/R}
 \,\dd\theta.
 }
 \label{R:cf:eq:matrixinversion}
\end{equation}
\end{theorem}

\begin{proof}
On $E_{k,\sigma}$,
$e^{i\theta\curl}$ is multiplication by
$e^{i\theta\mu_{k,\sigma}}$.
Thus \eqref{R:cf:eq:Fourier} follows from shell orthogonality.
Moreover
\[
 \sum_{k,\sigma}
 \|\cN_{k,\sigma}\|_1
 =
 \sum_{k,\sigma}
 \tr\cN_{k,\sigma}
 =
 \langle Av,v\rangle
 <\infty.
\]
The matrix Fourier series is therefore absolutely convergent in trace norm.
Fourier inversion on the circle of length $2\pi R$ gives
\eqref{R:cf:eq:matrixinversion}.
\end{proof}

\begin{theorem}[Operator-valued positive definiteness]
\label{R:cf:thm:positivefunction}
The periodic function $\Phi_v$ is positive definite: for arbitrary
$\theta_1,\ldots,\theta_m$ and
$X_1,\ldots,X_m\in\cK$,
\begin{equation}
 \boxed{
 \sum_{i,j=1}^m
 \left\langle
 X_i,
 \Phi_v(\theta_i-\theta_j)X_j
 \right\rangle_{\cK}
 \ge0.
 }
 \label{R:cf:eq:positivefunction}
\end{equation}
\end{theorem}

\begin{proof}
Use \eqref{R:cf:eq:Fourier}:
\[
\begin{aligned}
 &\sum_{i,j}
 \langle
 X_i,\Phi(\theta_i-\theta_j)X_j
 \rangle\\
 &\quad=
 \sum_\alpha
 \left\langle
 \sum_i e^{-i\mu_\alpha\theta_i}X_i,
 \cN_\alpha
 \sum_j e^{-i\mu_\alpha\theta_j}X_j
 \right\rangle
 \ge0.
\end{aligned}
\]
\end{proof}

This is the concrete Navier--Stokes realization of an operator-valued
Bochner/Herglotz representation \cite{Minh2016}.

\subsection{Infinite-band signed-shell energy tomography}

Take the trace:
\begin{equation}
 \phi_v(\theta)
 :=
 \tr\Phi_v(\theta).
 \label{R:cf:eq:phi}
\end{equation}
By \cref{R:cf:thm:momentrepresentation,R:cf:thm:Fourier},
\begin{equation}
 \boxed{
 \phi_v(\theta)
 =
 \sum_{k\ge1,\sigma=\pm}
 \kappa_k
 \|v_{k,\sigma}\|_2^2
 e^{i\sigma(k+2)\theta/R}.
 }
 \label{R:cf:eq:scalarFourier}
\end{equation}

\begin{theorem}[Exact infinite-band energy tomography]
\label{R:cf:thm:infiniteTomography}
Let $v$ satisfy $\langle Av,v\rangle<\infty$.  Then every signed-shell
energy is determined by the single scalar periodic function
$\phi_v$:
\begin{equation}
 \boxed{
 \|v_{k,\sigma}\|_2^2
 =
 \frac{1}{\kappa_k}
 \frac{1}{2\pi R}
 \int_0^{2\pi R}
 \phi_v(\theta)
 e^{-i\sigma(k+2)\theta/R}
 \,\dd\theta.
 }
 \label{R:cf:eq:energyinversion}
\end{equation}
No finite-band hypothesis is required.
\end{theorem}

\begin{proof}
The Fourier coefficient at
$n=\sigma(k+2)$ in \eqref{R:cf:eq:scalarFourier} is
$\kappa_k\|v_{k,\sigma}\|_2^2$.  Since $k\ge1$, $\kappa_k>0$.
\end{proof}

\begin{remark}[Forbidden low harmonics]
For $v\perp\cK$, the Fourier coefficients of $\phi_v$ and $\Phi_v$ vanish
at $n=0,\pm1,\pm2$.  The response circle therefore carries a built-in
spectral gap between the zero harmonic and the first high-mode frequencies
$n=\pm3$.
\end{remark}

Normalize by the total dissipation weight:
\begin{equation}
 \psi_v(\theta)
 :=
 \frac{\phi_v(\theta)}{\phi_v(0)}.
 \label{R:cf:eq:psi}
\end{equation}
Then $\psi_v$ is the characteristic function of the probability distribution
\begin{equation}
 p_{k,\sigma}
 =
 \frac{
 \kappa_k\|v_{k,\sigma}\|_2^2
 }{
 \langle Av,v\rangle
 }.
 \label{R:cf:eq:pk}
\end{equation}

\begin{corollary}[Dissipation-weighted curl statistics]
\label{R:cf:cor:curlstats}
If the required moments are finite, then
\begin{align}
 \frac1{i}\psi_v'(0)
 &=
 \sum_{k,\sigma}
 p_{k,\sigma}\mu_{k,\sigma},
 \label{R:cf:eq:meanmu}\\
 -\psi_v''(0)
 -
 \left|\psi_v'(0)\right|^2
 &=
 \operatorname{Var}_{p}(\mu)
 \ge0.
 \label{R:cf:eq:varmu}
\end{align}
The variance vanishes if and only if the field is supported on one signed
curl shell.
\end{corollary}

\subsection{Stokes dynamics as heat flow on the response circle}

Consider
\begin{equation}
       \partial_tv+\nu Av=0.
       \label{R:cf:eq:Stokes}
\end{equation}
Then
\[
       v_{k,\sigma}(t)
       =
       e^{-\nu\kappa_kt}v_{k,\sigma}(0)
\]
and hence
\begin{equation}
 \cN_{k,\sigma}(t)
 =
 e^{-2\nu\kappa_kt}
 \cN_{k,\sigma}(0).
 \label{R:cf:eq:Nt}
\end{equation}

\begin{theorem}[Matrix heat equation on the response circle]
\label{R:cf:thm:heatcircle}
Under \eqref{R:cf:eq:Stokes},
\begin{equation}
 \boxed{
 \partial_t\Phi
 =
 2\nu
 \left(
 \partial_\theta^2+\frac4{R^2}
 \right)\Phi
 }
 \label{R:cf:eq:heatcircle}
\end{equation}
on the circle
$\mathbb R/(2\pi R\mathbb Z)$.
The same equation holds for
$\phi=\tr\Phi$.
\end{theorem}

\begin{proof}
For one Fourier mode,
\[
 2\nu
 \left(
 \partial_\theta^2+\frac4{R^2}
 \right)
 e^{i\mu\theta}
 =
 -2\nu
 \left(
 \mu^2-\frac4{R^2}
 \right)
 e^{i\mu\theta}
 =
 -2\nu\kappa
 e^{i\mu\theta}.
\]
Compare with \eqref{R:cf:eq:Nt} and sum the absolutely convergent Fourier series.
\end{proof}

\begin{remark}
The auxiliary circle is not physical space.  It is the Pontryagin dual of
the integer-spaced curl spectrum.  The heat equation
\eqref{R:cf:eq:heatcircle} is therefore an exact spectral representation of
viscous relaxation, not an additional spatial dimension.
\end{remark}

\begin{corollary}[Monotone Fourier coefficient decay]
Every matrix Fourier coefficient evolves independently:
\[
 \widehat\Phi_{\sigma(k+2)}(t)
 =
 e^{-2\nu\kappa_kt}
 \widehat\Phi_{\sigma(k+2)}(0).
\]
Thus the response circle diagonalizes the complete matrix-valued Stokes
relaxation.
\end{corollary}

\subsection{What matrix tomography reconstructs}

The Fourier inversion theorem recovers
$\cN_{k,\sigma}$, not necessarily the full vector
$v_{k,\sigma}$.

\begin{proposition}[High-shell noninjectivity]
\label{R:cf:prop:noninjective}
Fix $k\ge3$ and $\sigma$.  Consider the normalized quadratic response map
\begin{equation}
 \mathcal R_{k,\sigma}:
 \mathbb RP^{d_k-1}
 \longrightarrow
 \operatorname{Sym}_6,
 \qquad
 [v]
 \longmapsto
 \frac{\cN_{k,\sigma}(v)}{\|v\|_2^2}.
 \label{R:cf:eq:responsemap}
\end{equation}
This map cannot be injective.
\end{proposition}

\begin{proof}
The domain is a compact manifold of real dimension
\[
       d_k-1=(k+1)(k+3)-1.
\]
For $k\ge3$ this is at least $23$.  The real vector space
$\operatorname{Sym}_6$ has dimension $21$.  A compact manifold of topological
dimension greater than $21$ cannot embed continuously in $\mathbb R^{21}$.
Since \eqref{R:cf:eq:responsemap} is continuous, it is not injective.
\end{proof}

\begin{remark}
Thus even complete knowledge of the response circle recovers the full
\emph{response geometry} of every signed shell and its energy, but not the
internal phase state of sufficiently high shells.  Additional covariants
would be required for state-level orbit reconstruction.
\end{remark}

\subsection{Information hierarchy, limitations, and relation to prior theory}

The combined structure separates four inverse problems.

\begin{enumerate}[label=(\roman*)]
\item
\textbf{State reconstruction.}
The response fields form a lossless transform; two fixed same-chirality
channels suffice globally.

\item
\textbf{Gram geometry.}
The \(6\times6\) response tensor retains dissipation, chiral traces and
isometry covariants, but it can lose internal shell phase information.

\item
\textbf{Moment tomography.}
The Stokes-weighted hierarchy recovers spectral moments and finite-band
energy distributions through Hankel/Vandermonde reconstruction.

\item
\textbf{Fourier tomography.}
The curl-response characteristic function recovers every signed-shell
response tensor and every signed-shell energy without a finite-band
hypothesis.
\end{enumerate}

General \(g\)-frames and operator-valued frames were developed independently
of fluid mechanics \cite{Sun2006,KaftalLarsonZhang2009}.  Matrix-valued
moment theory and operator-valued Bochner theorems provide abstract
reconstruction frameworks, while the present setting supplies a concrete
hydrodynamic measure with support fixed by the round-sphere curl spectrum.
The results should therefore be read as an exact geometric realization, not
as a claim that frame or moment theory itself is new.

Three limitations are essential.  First, Gram compression need not determine
a high-dimensional shell state.  Second, exact moment inversion from finitely
many moments requires finite spectral support or additional determinacy
information; the Fourier-circle theorem avoids that issue for
finite-dissipation states.  Third, none of the global \(L^2\) response
identities is an endpoint \(L^\infty\) continuation criterion for the
three-dimensional Navier--Stokes problem.



\section{Observable leakage, response-circle generators, and spectral selection}
\label{N:sec:synthesis}

The reduction and response theories developed above admit several further
consequences once they are used together.  This section records those
consequences separately because they connect objects that were introduced
for different purposes: the Clebsch--Gordan support of nonlinear leakage,
minimal Killing observability, the response characteristic function, and the
normalized spectral energy distribution.

\subsection{Clebsch--Gordan locality and parity of nonlinear interactions}

The complexified signed shell has the representation
\begin{equation}
 E_{k,\sigma}^{\mathbb C}
 \simeq
 V_{j_+}\boxtimes V_{j_-},
 \qquad
 j_\tau=\frac{k+1+\tau\sigma}{2}.
 \label{N:eq:shellrep}
\end{equation}
Since the Levi-Civita connection and Leray projection are isometry
equivariant, \(Q\) is an \(SO(4)\)-equivariant bilinear map.  Ref.~\cite{BateninPreeceKB2026} explicitly left the resulting
Clebsch--Gordan shell-selection question open; the next theorem supplies
that analysis.

\begin{theorem}[Interaction parity and spectral locality]
\label{N:thm:selection}
Suppose
\[
 \Pi_{\ell,\tau}
 Q(E_{k_1,\sigma_1},E_{k_2,\sigma_2})
 \neq0.
\]
Then
\begin{equation}
 \boxed{
 \ell\equiv k_1+k_2\pmod 2,
 \qquad
 \ell\le k_1+k_2.
 }
 \label{N:eq:paritylocality}
\end{equation}
More precisely,
\begin{equation}
 \boxed{
 \ell
 \le
 k_1+k_2+1-
 \left|\sigma_1+\sigma_2-\tau\right|.
 }
 \label{N:eq:refinedlocality}
\end{equation}
Thus, when
\(\sigma_1=\sigma_2=-\tau\),
one has the stronger bound
\[
       \ell\le k_1+k_2-2.
\]
\end{theorem}

\begin{proof}
A nonzero projected interaction produces a nonzero
\(SO(4)\)-intertwiner from
\(E_{k_1,\sigma_1}^{\mathbb C}\otimes
 E_{k_2,\sigma_2}^{\mathbb C}\)
to \(E_{\ell,\tau}^{\mathbb C}\).
For each \(SU(2)\) factor, the Clebsch--Gordan rule therefore requires
\[
 j_\pm(\ell,\tau)
 \le
 j_\pm(k_1,\sigma_1)+j_\pm(k_2,\sigma_2),
\]
and the output spin has the same integer/half-integer parity as the sum of
the two input spins.  Doubling the \(+\) spin gives
\[
 \ell+1+\tau
 \le
 k_1+k_2+2+\sigma_1+\sigma_2,
\]
while doubling the \(-\) spin gives
\[
 \ell+1-\tau
 \le
 k_1+k_2+2-\sigma_1-\sigma_2.
\]
Taking the smaller upper bound yields
\eqref{N:eq:refinedlocality}.  Modulo two,
\(2j_\pm(k,\sigma)\equiv k\), so the Clebsch--Gordan parity condition gives
\(\ell\equiv k_1+k_2\pmod2\).
\end{proof}

\begin{corollary}[Finite output band]
\label{N:cor:finiteout}
If a finite-dimensional candidate space \(W\) is supported in shell levels
\(k\le K\), then
\[
 Q(W,W)
 \subset
 \bigoplus_{\ell\le2K,\ \tau=\pm}E_{\ell,\tau}.
\]
In particular, every leakage field from \(W\) has Stokes eigenvalue at most
\[
 \kappa_{2K}
 =
 \frac{4K(K+2)}{R^2}.
\]
\end{corollary}

\subsection{Arbitrary two-probe observability}

The sharp theorem above used an orthonormal pair of Killing probes.  The
orthogonality assumption is not needed for injectivity.

\begin{theorem}[Angle-stable two-probe observability]
\label{N:thm:angledpair}
Let \(K,L\in\mathfrak K_\tau\) be linearly independent and let
\[
 G(K,L)
 =
 \begin{pmatrix}
 \langle K,K\rangle&\langle K,L\rangle\\
 \langle L,K\rangle&\langle L,L\rangle
 \end{pmatrix}
\]
be their \(L^2\) Gram matrix.  Then every
\(v\in D(A^{1/2})\cap\mathfrak K^\perp\) satisfies
\begin{equation}
 \boxed{
 2V_R
 \left(
 \|Q(K,v)\|_2^2+\|Q(L,v)\|_2^2
 \right)
 \ge
 \frac{\lambda_{\min}(G(K,L))}{3R^2}
 \|v\|_2^2.
 }
 \label{N:eq:angledpair}
\end{equation}
Hence any two linearly independent same-chirality probes determine the
complete non-Killing state.

If \(K,L\) are unit probes and \(\theta\) is their angle, then
\begin{equation}
 2V_R
 \left(
 \|Q(K,v)\|_2^2+\|Q(L,v)\|_2^2
 \right)
 \ge
 \frac{1-|\cos\theta|}{3R^2}\|v\|_2^2.
 \label{N:eq:anglebound}
\end{equation}
\end{theorem}

\begin{proof}
Choose an orthonormal basis \(e_1,e_2\) of
\(\operatorname{span}\{K,L\}\), and let \(C\) be the \(2\times2\) matrix of
the coordinates of \(K,L\) in that basis.  By linearity in the Killing
argument,
\[
 \binom{Q(K,v)}{Q(L,v)}
 =
 C
 \binom{Q(e_1,v)}{Q(e_2,v)}.
\]
Therefore
\[
\begin{aligned}
 &\|Q(K,v)\|_2^2+\|Q(L,v)\|_2^2\\
 &\qquad\ge
 \lambda_{\min}(C^*C)
 \left(
 \|Q(e_1,v)\|_2^2+\|Q(e_2,v)\|_2^2
 \right).
\end{aligned}
\]
The nonzero eigenvalues of \(C^*C\) are those of
\(CC^*=G(K,L)\).  Apply the sharp orthonormal two-probe bound
\eqref{R:eq:twobound}.  For unit probes the Gram eigenvalues are
\(1\pm|\cos\theta|\).
\end{proof}

There is also a uniform one-probe result on the odd shell sector.

\begin{theorem}[Uniform one-probe coercivity on odd levels]
\label{N:thm:oddcoercive}
Let \(K\in\mathfrak K_\tau\) be \(L^2\)-normalized and let
\[
 \mathcal H_{\mathrm{odd}}
 =
 \widehat{\bigoplus}_{\substack{k\ge1\ {\rm odd}\\\sigma=\pm}}
 E_{k,\sigma}.
\]
Then, for
\(v\in D(A^{1/2})\cap\mathcal H_{\mathrm{odd}}\),
\begin{equation}
 \boxed{
 2V_R\|Q(K,v)\|_2^2
 \ge
 \frac{1}{18R^2}\|v\|_2^2.
 }
 \label{N:eq:oddcoercive}
\end{equation}
The constant \(1/(18R^2)\) is sharp.
\end{theorem}

\begin{proof}
On \(E_{k,\sigma}\), align \(K\) with a Cartan generator of the
\(SU(2)_\tau\) factor.  The chiral three-axis response has scalar eigenvalue
\[
 d_\tau(k,\sigma)
 =
 \frac{\kappa_k}{2}
 \left(1-\frac{2\tau\sigma}{k+2}\right).
\]
For spin \(j_\tau=(k+1+\tau\sigma)/2\), a single axis has squared response
weights proportional to
\(m^2/[j_\tau(j_\tau+1)]\).  When \(k\) is odd,
\(j_\tau\) is half-integer and \(|m|\ge1/2\).  Hence the minimum one-axis
eigenvalue is
\[
 \frac{d_\tau(k,\sigma)}
      {4j_\tau(j_\tau+1)}
 =
 \begin{cases}
 \dfrac{k^2}{2R^2(k+2)^2},&\tau\sigma=+1,\\[2mm]
 \dfrac{(k+4)^2}{2R^2(k+2)^2},&\tau\sigma=-1.
 \end{cases}
\]
Over odd \(k\ge1\) and both signs, the minimum is
\(1/(18R^2)\), attained at \(k=1\) and \(\tau\sigma=+1\).
Shell orthogonality gives \eqref{N:eq:oddcoercive}.
\end{proof}

\subsection{Two-probe detection of the full leakage variety}

Let \(W\subset\mathfrak K^\perp\) be a finite-dimensional Stokes-invariant
candidate reduction and define its bilinear leakage map
\begin{equation}
 \mathcal L_W(u,v)
 :=
 (I-P_W)Q(u,v),
 \qquad u,v\in W.
 \label{N:eq:leakmap}
\end{equation}
For an orthonormal basis \(\{e_i\}\) of \(W\), set
\[
 \Lambda(W)
 =
 \sum_{i,j}\|\mathcal L_W(e_i,e_j)\|_2^2.
\]
The symmetric projected convection has no Killing component:
for every Killing field \(K_0\),
\begin{equation}
 \langle Q(u,v),K_0\rangle=0.
 \label{N:eq:noKillingQ}
\end{equation}
Indeed, integration by parts followed by the Killing equation cancels the
two polarized terms.  Thus every leakage field lies in
\(\mathfrak K^\perp\).

\begin{theorem}[Two-probe leakage certificate]
\label{N:thm:leakcertificate}
Let \(K,L\in\mathfrak K_\tau\) be linearly independent and define
\begin{align}
 \Lambda_{K,L}(W)
 &:=
 2V_R\sum_{i,j}
 \Bigl(
 \|Q(K,\mathcal L_W(e_i,e_j))\|_2^2
 \notag\\
 &\hspace{7em}
 +\|Q(L,\mathcal L_W(e_i,e_j))\|_2^2
 \Bigr).
 \label{N:eq:observedleak}
\end{align}
If \(W\) is supported in levels \(k\le K_{\max}\), then
\begin{equation}
\begin{aligned}
 \frac{\lambda_{\min}(G(K,L))}{3R^2}\Lambda(W)
 &\le
 \Lambda_{K,L}(W)\\
 &\le
 \lambda_{\max}(G(K,L))
 \kappa_{2K_{\max}}\Lambda(W).
\end{aligned}
\label{N:eq:leakequiv}
\end{equation}
Consequently,
\begin{equation}
 \boxed{
 \Lambda(W)=0
 \quad\Longleftrightarrow\quad
 \Lambda_{K,L}(W)=0.
 }
 \label{N:eq:leakzeroequiv}
\end{equation}

Moreover, if \(\Phi\) is the algebraic leakage section on a product of
Grassmannians and
\[
 \Psi_{K,L}
 :=
 \bigl(Q(K,\cdot),Q(L,\cdot)\bigr)\circ\Phi,
\]
then at every zero-leakage point
\begin{equation}
 \boxed{
 \ker d\Psi_{K,L}=\ker d\Phi,
 \qquad
 \operatorname{rank}d\Psi_{K,L}
 =
 \operatorname{rank}d\Phi.
 }
 \label{N:eq:tangentvisible}
\end{equation}
Thus two Killing-response channels detect not only exact closure but also
the full infinitesimal deformation geometry of the leakage variety.
\end{theorem}

\begin{proof}
Apply \cref{N:thm:angledpair} to each leakage field and sum over \(i,j\).
This gives the lower bound.  For the upper bound, choose an orthonormal basis
of \(\operatorname{span}\{K,L\}\).  The largest eigenvalue of \(G(K,L)\)
bounds the pair response by that orthonormal pair, and the latter is bounded
by the full six-channel Casimir sum.  Hence
\[
\begin{aligned}
 \Lambda_{K,L}(W)
 &\le
 \lambda_{\max}(G(K,L))
 \sum_{i,j}
 \bigl\langle
 A\mathcal L_W(e_i,e_j),\\
 &\hspace{9em}
 \mathcal L_W(e_i,e_j)
 \bigr\rangle.
\end{aligned}
\]
By \cref{N:thm:selection,N:cor:finiteout}, every leakage field is supported
in levels at most \(2K_{\max}\), yielding the upper estimate.

The response pair is injective on \(\mathfrak K^\perp\), so composing the
leakage section with it does not change the zero set.  At a zero,
\(d\Psi_{K,L}\) is the fixed injective response pair composed with \(d\Phi\).
An injective postcomposition preserves both kernel and rank.
\end{proof}

\begin{corollary}[One-probe closure test for opposite shell parity]
\label{N:cor:oneprobeclosure}
Let
\[
 W=L_\alpha\oplus L_\beta,
 \qquad
 L_\alpha\subset E_{k_1,\sigma_1},
 \quad
 L_\beta\subset E_{k_2,\sigma_2},
\]
with \(k_1+k_2\) odd.  For any \(L^2\)-normalized nonzero
\(K\in\mathfrak K_\tau\), define
\[
 \Lambda_K(W)
 =
 2V_R
 \|Q(K,\mathcal L_W)\|_{\HS}^2.
\]
Then
\begin{equation}
 \boxed{
 \frac{1}{18R^2}\Lambda(W)
 \le
 \Lambda_K(W)
 \le
 \kappa_{k_1+k_2}\Lambda(W),
 }
 \label{N:eq:oneprobeclosure}
\end{equation}
and hence
\[
       W\text{ is exact}
       \quad\Longleftrightarrow\quad
       \Lambda_K(W)=0.
\]
\end{corollary}

\begin{proof}
Single-shell terms vanish identically, so only the cross interaction
\(Q(L_\alpha,L_\beta)\) remains.  By
\cref{N:thm:selection}, every output shell has odd index because
\(k_1+k_2\) is odd, and its level is at most \(k_1+k_2\).
Apply \cref{N:thm:oddcoercive} and the Casimir upper bound.
\end{proof}

\subsection{The response circle as a master generating object}

Let
\[
 \mathscr D_\theta=-i\partial_\theta,
 \qquad
 \mathscr L_\theta
 =
 \mathscr D_\theta^2-\frac4{R^2}
 =
 -\partial_\theta^2-\frac4{R^2}
\]
on the circle of length \(2\pi R\), restricted to the high harmonics
\(|n|\ge3\).  On the Fourier character
\(e^{i\mu_{k,\sigma}\theta}\),
\[
 \mathscr D_\theta=\mu_{k,\sigma},
 \qquad
 \mathscr L_\theta=\kappa_k.
\]

\begin{theorem}[Master response-circle functional calculus]
\label{N:thm:mastercircle}
Let \(\Phi_v\) be the operator-valued response characteristic function from
\cref{R:cf:thm:Fourier}.  Whenever the corresponding moments are finite,
\begin{align}
 \boxed{
 \cC_n(v)
 =
 \left.
 \mathscr D_\theta^n\Phi_v(\theta)
 \right|_{\theta=0},
 }
 \label{N:eq:curljet}\\
 \boxed{
 \cM_v^{[s]}
 =
 \left.
 \mathscr L_\theta^s\Phi_v(\theta)
 \right|_{\theta=0}.
 }
 \label{N:eq:stokesjet}
\end{align}
For noninteger \(s\), \(\mathscr L_\theta^s\) is the positive Fourier
multiplier on the high-harmonic subspace.

Define the energy-response characteristic function
\begin{equation}
 \mathscr E_v(\theta)
 :=
 \mathscr L_\theta^{-1}\Phi_v(\theta),
 \qquad
 \varepsilon_v(\theta)
 :=
 \tr\mathscr E_v(\theta).
 \label{N:eq:energycircle}
\end{equation}
Then
\begin{equation}
 \boxed{
 \varepsilon_v(\theta)
 =
 \sum_{k\ge1,\sigma=\pm}
 \|v_{k,\sigma}\|_2^2
 e^{i\sigma(k+2)\theta/R}.
 }
 \label{N:eq:energyFourier}
\end{equation}
Consequently,
\begin{align}
 \|v\|_2^2
 &=
 \varepsilon_v(0),
 \label{N:eq:energyat0}\\
 \langle \curl^r v,v\rangle
 &=
 \left.
 \mathscr D_\theta^r\varepsilon_v
 \right|_{\theta=0},
 \label{N:eq:curlmomenergy}\\
 \langle A^n v,v\rangle
 &=
 \left.
 \mathscr L_\theta^n\varepsilon_v
 \right|_{\theta=0},
 \label{N:eq:stokemomenergy}\\
 \langle A^n\curl^{-1}v,v\rangle
 &=
 \left.
 \mathscr D_\theta^{-1}
 \mathscr L_\theta^n\varepsilon_v
 \right|_{\theta=0}.
 \label{N:eq:signedmomenergy}
\end{align}
Thus the Stokes hierarchy, curl hierarchy, kinetic energy, helicity and
signed inverse-curl moments are all jets or Fourier multipliers of one
periodic function.
\end{theorem}

\begin{proof}
The Fourier representation
\[
 \Phi_v(\theta)
 =
 \sum_{k,\sigma}
 e^{i\mu_{k,\sigma}\theta}\cN_{k,\sigma}
\]
and the identities
\[
 \mathscr D_\theta e^{i\mu\theta}
 =
 \mu e^{i\mu\theta},
 \qquad
 \mathscr L_\theta e^{i\mu\theta}
 =
 \left(\mu^2-\frac4{R^2}\right)e^{i\mu\theta}
 =
 \kappa e^{i\mu\theta}
\]
give \eqref{N:eq:curljet}--\eqref{N:eq:stokesjet}.
Since
\(\tr\cN_{k,\sigma}=\kappa_k\|v_{k,\sigma}\|_2^2\),
applying \(\mathscr L_\theta^{-1}\) and taking the trace gives
\eqref{N:eq:energyFourier}.  The remaining identities are its Fourier
moments.
\end{proof}

\subsection{Extreme spectral characters and exact same-shell sectors}

For \(v\neq0\), normalize
\begin{align}
 \chi_v(\theta)
 &:=
 \frac{\varepsilon_v(\theta)}{\varepsilon_v(0)}
 =
 \sum_{k,\sigma}
 p_{k,\sigma}
 e^{i\sigma(k+2)\theta/R},
 \label{N:eq:normalizedenergycircle}\\
 p_{k,\sigma}
 &:=
 \frac{\|v_{k,\sigma}\|_2^2}{\|v\|_2^2}.
 \label{N:eq:spectralprob}
\end{align}
This is a positive-definite function and the characteristic function of a
probability measure on the allowed signed curl lattice.

\begin{theorem}[Spectral extremality of same-shell exact states]
\label{N:thm:extreme}
Within the convex set of normalized energy-response functions
\eqref{N:eq:normalizedenergycircle}, \(\chi_v\) is an extreme point if and
only if \(v\) is supported on one signed curl shell.  In that case
\begin{equation}
 \boxed{
 \chi_v(\theta)
 =
 e^{i\sigma(k+2)\theta/R},
 }
 \label{N:eq:character}
\end{equation}
and \(Q(v,v)=0\).

Thus every Grassmannian family
\(\Gr(d,E_{k,\sigma})\) of exact same-shell reductions is collapsed by
scalar spectral tomography to the single extreme character
\eqref{N:eq:character}.
\end{theorem}

\begin{proof}
The Fourier coefficients in
\eqref{N:eq:normalizedenergycircle} form a probability vector.
The extreme points of a probability simplex are the Dirac masses, hence the
extreme positive-definite functions are the single characters.  A Dirac
spectral measure is equivalent to support in one signed shell.  The
single-shell cancellation then gives \(Q(v,v)=0\).
\end{proof}

\subsection{Fisher--Rao Stokes selection}

The normalized shell energies reveal a second geometry.  General
Fisher--Rao gradient-flow theory is independent of the present fluid
construction \cite{Carrillo2026}; here it arises exactly from Stokes
spectral decay.

\begin{theorem}[Fisher--Rao spectral selection law]
\label{N:thm:fisher}
Consider the linear Stokes evolution
\[
       \partial_t v+\nu Av=0
\]
and set
\[
\begin{aligned}
 E_\alpha(t)&=\|v_\alpha(t)\|_2^2,\\
 p_\alpha(t)&=
 \frac{E_\alpha(t)}{\sum_\beta E_\beta(t)},\\
 \bar\kappa(t)&=
 \sum_\alpha p_\alpha(t)\kappa_\alpha.
\end{aligned}
\]
where \(\alpha=(k,\sigma)\).
Then
\begin{equation}
 \boxed{
 \dot p_\alpha
 =
 -2\nu
 \bigl(\kappa_\alpha-\bar\kappa\bigr)p_\alpha.
 }
 \label{N:eq:replicator}
\end{equation}
On every finite spectral simplex, if
\[
 g_p(\xi,\eta)
 =
 \sum_\alpha\frac{\xi_\alpha\eta_\alpha}{p_\alpha}
\]
is the Fisher--Rao metric and
\(\mathcal K(p)=\sum_\alpha\kappa_\alpha p_\alpha\), then
\begin{equation}
 \boxed{
 \dot p=-2\nu\,\operatorname{grad}_{FR}\mathcal K.
 }
 \label{N:eq:FRgradient}
\end{equation}
Moreover,
\begin{equation}
 \boxed{
 \frac{\dd}{\dd t}\bar\kappa
 =
 -2\nu\operatorname{Var}_{p}(\kappa)
 \le0.
 }
 \label{N:eq:meanKdecay}
\end{equation}
Equality holds exactly when the normalized energy is supported on one
Stokes level.
\end{theorem}

\begin{proof}
The shell energies satisfy
\[
 E_\alpha(t)
 =
 E_\alpha(0)e^{-2\nu\kappa_\alpha t}.
\]
Differentiating their normalization gives
\eqref{N:eq:replicator}.  For a tangent vector \(\xi\) with
\(\sum_\alpha\xi_\alpha=0\),
\[
 d\mathcal K[\xi]
 =
 \sum_\alpha\kappa_\alpha\xi_\alpha
 =
 g_p\!\left(
 \bigl(p_\alpha(\kappa_\alpha-\bar\kappa)\bigr)_\alpha,
 \xi
 \right),
\]
which proves \eqref{N:eq:FRgradient}.  Finally,
\[
 \dot{\bar\kappa}
 =
 -2\nu
 \sum_\alpha
 p_\alpha(\kappa_\alpha-\bar\kappa)\kappa_\alpha
 =
 -2\nu\operatorname{Var}_{p}(\kappa).
\]
\end{proof}

\begin{corollary}[Exact viscous spectral selection]
\label{N:cor:selection}
Let \(k_*\) be the smallest active Stokes level and let
\[
 P_*(0)
 =
 p_{k_*,+}(0)+p_{k_*,-}(0)>0.
\]
If \(\Delta\kappa\) is the gap from \(\kappa_{k_*}\) to the next active
Stokes eigenvalue, then
\begin{equation}
 1-P_*(t)
 \le
 \frac{1-P_*(0)}{P_*(0)}
 e^{-2\nu\Delta\kappa t}.
 \label{N:eq:selectionrate}
\end{equation}
The ratio
\(p_{k_*,+}(t):p_{k_*,-}(t)\)
is constant.  Thus viscosity selects the lowest active Stokes
\emph{level}, but does not by itself equalize or select its two chiralities.
\end{corollary}

\subsection{Nonlinear response-circle transfer}

The same normalized spectral picture survives the full nonlinear equation as
an exact decomposition into a viscous gradient term and an energy-conserving
transfer term.

For a sufficiently regular unforced high-mode solution
\begin{equation}
 \partial_t v+\nu Av+Q(v,v)+D(m)v=0,
 \label{N:eq:highmode}
\end{equation}
define
\begin{equation}
 \mathcal T_{k,\sigma}
 :=
 -2
 \left\langle
 \Pi_{k,\sigma}Q(v,v),
 v_{k,\sigma}
 \right\rangle
 \label{N:eq:shelltransfer}
\end{equation}
and its response-circle transform
\begin{equation}
 \mathcal T_v(\theta)
 :=
 \sum_{k,\sigma}
 \mathcal T_{k,\sigma}
 e^{i\sigma(k+2)\theta/R}.
 \label{N:eq:transfercircle}
\end{equation}

\begin{theorem}[Nonlinear response-circle balance]
\label{N:thm:nonlinearcircle}
The energy-response characteristic function
\(\varepsilon_v\) satisfies
\begin{equation}
 \boxed{
 \partial_t\varepsilon_v
 =
 2\nu
 \left(
 \partial_\theta^2+\frac4{R^2}
 \right)
 \varepsilon_v
 +
 \mathcal T_v.
 }
 \label{N:eq:nonlinearcircle}
\end{equation}
The nonlinear transfer source obeys the two exact conservation constraints
\begin{equation}
 \boxed{
 \mathcal T_v(0)=0,
 \qquad
 \left.
 \mathscr D_\theta\mathcal T_v
 \right|_{\theta=0}=0.
 }
 \label{N:eq:doublezero}
\end{equation}
They are respectively the energy and helicity conservation laws of the
Euler nonlinearity.

If
\(E(t)=\|v(t)\|_2^2\) and
\(p_\alpha=E_\alpha/E\), then
\begin{equation}
 \boxed{
 \dot p_\alpha
 =
 -2\nu(\kappa_\alpha-\bar\kappa)p_\alpha
 +
 \frac{\mathcal T_\alpha}{E}.
 }
 \label{N:eq:gradienttransfer}
\end{equation}
Thus normalized spectral dynamics decomposes exactly into a Fisher--Rao
viscous gradient flow and a zero-sum nonlinear redistribution.  In
particular,
\begin{equation}
 \boxed{
 \dot{\bar\kappa}
 =
 -2\nu\operatorname{Var}_{p}(\kappa)
 +
 \frac1E
 \sum_\alpha\kappa_\alpha\mathcal T_\alpha.
 }
 \label{N:eq:meanKnonlinear}
\end{equation}
\end{theorem}

\begin{proof}
The Killing drift preserves each signed shell and is skew-adjoint there, so
it contributes nothing to the shell energies.  Taking the \(L^2\) inner
product of \eqref{N:eq:highmode} with \(v_{k,\sigma}\) gives
\[
 \dot E_{k,\sigma}
 =
 -2\nu\kappa_kE_{k,\sigma}
 +
 \mathcal T_{k,\sigma}.
\]
Fourier summation and
\[
 2\nu
 \left(\partial_\theta^2+\frac4{R^2}\right)
 e^{i\mu_{k,\sigma}\theta}
 =
 -2\nu\kappa_k
 e^{i\mu_{k,\sigma}\theta}
\]
give \eqref{N:eq:nonlinearcircle}.

The zeroth transfer moment is
\[
 \sum_{k,\sigma}\mathcal T_{k,\sigma}
 =
 -2\langle Q(v,v),v\rangle
 =
 0.
\]
The first signed-curl moment is
\[
 \sum_{k,\sigma}
 \mu_{k,\sigma}\mathcal T_{k,\sigma}
 =
 -2\langle Q(v,v),\curl v\rangle
 =
 0,
\]
because
\(\nabla_vv=\gradg(|v|^2/2)-v\times\curl v\):
the gradient is orthogonal to \(\curl v\) and
\((v\times\curl v)\cdot\curl v=0\).
This proves \eqref{N:eq:doublezero}.

Normalize the shell-energy equation by \(E(t)\).  Since the nonlinear
transfer sums to zero,
\(\dot E=-2\nu\bar\kappa E\), which gives
\eqref{N:eq:gradienttransfer}.  Multiplication by \(\kappa_\alpha\) and
summation gives \eqref{N:eq:meanKnonlinear}.
\end{proof}

\begin{corollary}[Finite-band conservation factorization]
\label{N:cor:transferfactor}
If \(v\) has finite spectral support, write
\(z=e^{i\theta/R}\).
Then the Laurent polynomial
\[
 P_v(z)
 =
 \sum_{k,\sigma}
 \mathcal T_{k,\sigma}
 z^{\sigma(k+2)}
\]
has a double root at \(z=1\).  Hence
\begin{equation}
 \boxed{
 P_v(z)=(z-1)^2S_v(z)
 }
 \label{N:eq:transferfactor}
\end{equation}
for another Laurent polynomial \(S_v\).
\end{corollary}

\subsection{Quantitative information loss under Gram compression}

The two-probe transform is injective, whereas the pure-shell Gram tensor is
necessarily many-to-one at high shell number.  The chiral trace constraints
sharpen the dimension obstruction.

\begin{theorem}[Lower bound on Gram-compression fiber dimension]
\label{N:thm:gramfiber}
Fix \(k\ge1\) and \(\sigma\), and consider the normalized pure-shell Gram map
\[
 \Gamma_{k,\sigma}:
 \mathbb RP^{d_k-1}
 \longrightarrow
 \operatorname{Sym}_6,
 \qquad
 [v]\longmapsto
 \frac{\cM_v}{\|v\|_2^2}.
\]
Its image lies in the affine subspace
\begin{equation}
\begin{aligned}
 \mathcal A_{k,\sigma}
 =
 \bigl\{
 M\in\operatorname{Sym}_6:\;&
 \tr_{\mathfrak K_+}M=d_+^{(k,\sigma)},\\
 &
 \tr_{\mathfrak K_-}M=d_-^{(k,\sigma)}
 \bigr\}.
\end{aligned}
\label{N:eq:gramaffine}
\end{equation}
which has dimension \(19\).

On every constant-rank stratum of \(\Gamma_{k,\sigma}\), every local fiber
has dimension at least
\begin{equation}
 \boxed{
 d_k-20
 =
 k^2+4k-17.
 }
 \label{N:eq:fiberdim}
\end{equation}
In particular, for \(k\ge3\) Gram compression loses at least four continuous
projective degrees of freedom at every regular point of maximal differential
rank.
\end{theorem}

\begin{proof}
The pure-shell chiral partition fixes the two block traces independently
after normalization by \(\|v\|_2^2\).  Since
\(\dim\operatorname{Sym}_6=21\), the two independent affine trace
constraints reduce the ambient dimension to \(19\).

The domain has dimension \(d_k-1\).  On a constant-rank stratum with
differential rank \(r\le19\), the constant-rank theorem gives local fiber
dimension
\[
       d_k-1-r
       \ge
       d_k-20.
\]
For \(k=3\), this lower bound is \(24-20=4\), and it increases quadratically
with \(k\).
\end{proof}

\begin{corollary}[Exact information ladder on a pure shell]
\label{N:cor:informationladder}
For a fixed signed shell \(E_{k,\sigma}\):

\begin{enumerate}[label=(\roman*)]
\item scalar normalized spectral tomography is constant on
\(\mathbb RP^{d_k-1}\);

\item normalized Gram compression has local fibers of dimension at least
\(d_k-20\) on every constant-rank stratum;

\item any two linearly independent same-chirality raw Killing responses are
injective.

\end{enumerate}
Thus information loss occurs at compression, not at the Killing-response
analysis itself.
\end{corollary}


\section{Open problems}

The unified structure isolates several concrete problems.

\begin{enumerate}[label=(\arabic*)]
\item
Determine the principal orbit types and singular strata of
\(\Gr(d,E_{k,\sigma})/SO(4)\), and compute the first nontrivial irreducible
components of the multi-shell zero-leakage variety.

\item
Classify when the constant-drift Grassmannian geodesic condition is
compatible with nontrivial resonance-free shell splitting.

\item
Determine exact frame bounds for arbitrary nonorthogonal pairs of Killing
probes and the minimal spatial sampling needed to reconstruct their response
fields.

\item
Classify the image and orbit strata of the pure-shell Gram map
\(v\mapsto\mathscr M_v/\|v\|_2^2\), including extremal spectra and
cross-chiral ranks.

\item
Develop stable noisy versions of the finite-band Hankel reconstruction and
the infinite-band response-circle Fourier inversion.

\item
Use the nonlinear tensor balance and leakage geometry together to determine
when cross-shell interaction creates new response spectral rank.

\item
Find higher \(SO(4)\)-covariants that complement the Gram tensor and
separate generic isometry orbits inside high signed-curl shells.

\item
Study persistence of the reduction moduli, response frame, and response
circle under perturbations of the round metric.
\end{enumerate}

\section{Conclusion}

The Killing--Beltrami structure on the round three-sphere organizes both the
exact subsystems of Navier--Stokes flow and the information contained in
their symmetry responses.

On the reduction side, every subspace of a signed curl shell is an exact
pressure-linear Navier--Stokes system.  These systems form Grassmannian
moduli, acquire exact projector dynamics under Killing drift, and extend to
multi-shell reductions as the zero set of a real-algebraic leakage
functional.  Constant drift admits an exact geodesic criterion, an explicit
shell spectrum, finite resonance arithmetic, and special-orthogonal
transport.

On the reconstruction side, the six Killing responses factor the Stokes
operator and, after normalization, form a lossless Parseval transform.  Two
same-chirality probes already determine the complete non-Killing velocity,
while the one-probe problem has an exact shell-parity law.  Gram compression
produces a finite-dimensional response geometry; Stokes weighting produces
a positive moment hierarchy; curl weighting produces a positive
matrix-valued spectral measure whose periodic characteristic function
admits exact Fourier inversion.

The common chain is therefore
\[
\begin{aligned}
\text{signed-shell structure}
&\longrightarrow
\text{exact reductions}\\
&\longrightarrow
\text{Killing transport}\\
&\longrightarrow
\text{response tomography}.
\end{aligned}
\]
This is one geometric-spectral theory rather than two independent
manuscripts.


\section*{AI-Assisted Research Method}
OpenAI ChatGPT (GPT-5.6 Sol) was used as an AI-assisted research tool to
generate candidate algebraic deductions, test derivation paths, perform
algebraic consistency checks, and organize targeted literature searches.
The tool was directed through explicit mathematical prompts specifying the
relevant hypotheses, operators, and target identities. Candidate statements
were retained only when supported by explicit derivations in the manuscript
or by the cited sources; unsupported candidate claims were discarded. The
authors independently reviewed the retained AI-assisted mathematical claims
and source attributions and assume full responsibility for their accuracy
and interpretation.

\begin{acknowledgments}
In addition to the research-method use disclosed above, OpenAI ChatGPT
(GPT-5.6 Sol) assisted with manuscript organization and \LaTeX{}
preparation. The authors remain fully responsible for the accuracy,
interpretation, and scientific content of the manuscript.
\end{acknowledgments}

\section*{Data Availability}
No data were created or analyzed in this study.

\begin{thebibliography}{99}


\bibitem{BateninPreeceUnified2026}
B.~Batenin and A.~Preece,
``Geometric Reduction, Killing-Response Reconstruction, and Spectral
Tomography for Navier--Stokes Flow on the Round Three-Sphere,''
Zenodo, Version 1.0 (2026),
doi:10.5281/zenodo.22875853.


\bibitem{BateninPreeceKB2026}
B.~Batenin and A.~Preece,
``Killing--Beltrami Algebra and exact Reductions of Navier--Stokes Flow on
the Three-Sphere,''
Zenodo, Version 1.0 (2026),
doi:10.5281/zenodo.22874017.

\bibitem{BateninPreece2026}
B.~Batenin and A.~Preece,
``Finite-Dimensional Spectral Closure and a Relative Forced Singular Branch
for Navier--Stokes on the Round Three-Sphere,''
Zenodo, Version 1.0 (2026),
doi:10.5281/zenodo.22858542.

\bibitem{Bendokat2024}
T.~Bendokat, R.~Zimmermann, and P.-A. Absil,
``A Grassmann manifold handbook: basic geometry and computational aspects,''
Adv. Comput. Math. \textbf{50}, 6 (2024),
doi:10.1007/s10444-023-10090-8.

\bibitem{BoggarapuThangavelu2016}
P.~Boggarapu and S.~Thangavelu,
``Mixed norm estimates for the Riesz transforms on $SU(2)$,''
Publ. Mat. \textbf{60}, 171--190 (2016),
doi:10.5565/PUBLMAT\_60116\_06.

\bibitem{BogoyavlenskijFuchssteiner2005}
O.~I. Bogoyavlenskij and B.~Fuchssteiner,
``Exact NSE solutions with crystallographic symmetries and no transfer of
energy through the spectrum,''
J. Geom. Phys. \textbf{54}, 324--338 (2005),
doi:10.1016/j.geomphys.2004.10.003.

\bibitem{Buza2023}
G.~Buza,
``Spectral Submanifolds of the Navier--Stokes Equations,''
arXiv:2301.07898 (2023).


\bibitem{Carrillo2026}
J.~A. Carrillo, Y.~Chen, D.~Z. Huang, J.~Huang, and D.~Wei,
``Fisher--Rao Gradient Flow: Geodesic Convexity and Functional Inequalities,''
SIAM J. Math. Anal. \textbf{58}, 1062--1099 (2026),
doi:10.1137/25M176252X.

\bibitem{Capoferri2026}
M.~Capoferri and D.~Vassiliev,
``Beyond the Hodge theorem: Curl and asymmetric pseudodifferential projections,''
J. Lond. Math. Soc. \textbf{113}, e70431 (2026),
doi:10.1112/jlms.70431.

\bibitem{Edelman1998}
A.~Edelman, T.~A. Arias, and S.~T. Smith,
``The geometry of algorithms with orthogonality constraints,''
SIAM J. Matrix Anal. Appl. \textbf{20}, 303--353 (1998),
doi:10.1137/S0895479895290954.

\bibitem{EncisoPeraltaSalas2017}
A.~Enciso and D.~Peralta-Salas,
``Vortex reconnection in the three dimensional Navier--Stokes equations,''
Adv. Math. \textbf{309}, 452--486 (2017),
doi:10.1016/j.aim.2017.01.025.

\bibitem{EncisoReview2023}
A.~Enciso and D.~Peralta-Salas,
``Beltrami fields and knotted vortex structures in incompressible fluid
flows,''
Bull. Lond. Math. Soc. \textbf{55}, 2819--2845 (2023),
doi:10.1112/blms.12780.

\bibitem{Fasino1995}
D.~Fasino,
``Spectral properties of Hankel matrices and numerical solutions of finite
moment problems,''
J. Comput. Appl. Math. \textbf{65}, 145--155 (1995),
doi:10.1016/0377-0427(95)00108-5.

\bibitem{FrameTowne2023}
P.~Frame and A.~Towne,
``Space-time POD and the Hankel matrix,''
PLoS ONE \textbf{18}, e0289637 (2023),
doi:10.1371/journal.pone.0289637.

\bibitem{IkedaTaniguchi1978}
A.~Ikeda and Y.~Taniguchi,
``Spectra and eigenforms of the Laplacian on $S^n$ and $P^n(\mathbb C)$,''
Osaka J. Math. \textbf{15}, 515--546 (1978).

\bibitem{KaftalLarsonZhang2009}
V.~Kaftal, D.~R. Larson, and S.~Zhang,
``Operator-Valued Frames,''
Trans. Am. Math. Soc. \textbf{361}, 6349--6385 (2009),
doi:10.1090/S0002-9947-09-04915-0.

\bibitem{KimseyTrachana2022}
D.~P. Kimsey and M.~Trachana,
``On a Solution of the Multidimensional Truncated Matrix-Valued Moment
Problem,''
Milan J. Math. \textbf{90}, 17--101 (2022),
doi:10.1007/s00032-021-00346-7.

\bibitem{LiXu2025}
D.~Li and Y.~Xu,
``Phase and norm retrievable operator valued frames,''
Adv. Oper. Theory \textbf{10}, 82 (2025),
doi:10.1007/s43036-025-00449-7.

\bibitem{MahalovTitiLeibovich1990}
A.~Mahalov, E.~S. Titi, and S.~Leibovich,
``Invariant helical subspaces for the Navier--Stokes equations,''
Arch. Ration. Mech. Anal. \textbf{112}, 193--222 (1990),
doi:10.1007/BF00381234.

\bibitem{Minh2016}
H.~Q. Minh,
``Operator-Valued Bochner Theorem, Fourier Feature Maps for Operator-Valued
Kernels, and Vector-Valued Learning,''
arXiv:1608.05639 (2016).

\bibitem{SamavakiTuomela2020}
M.~Samavaki and J.~Tuomela,
``Navier--Stokes equations on Riemannian manifolds,''
J. Geom. Phys. \textbf{148}, 103543 (2020),
doi:10.1016/j.geomphys.2019.103543.

\bibitem{Sun2006}
W.~Sun,
``G-frames and g-Riesz bases,''
J. Math. Anal. Appl. \textbf{322}, 437--452 (2006),
doi:10.1016/j.jmaa.2005.09.039.

\bibitem{SunZalar2026}
S.~Sun and A.~Zalar,
``Matrix Fej\'er--Riesz type theorem for a union of an interval and a point,''
J. Pure Appl. Algebra \textbf{230}, 108173 (2026),
doi:10.1016/j.jpaa.2026.108173.

\end{thebibliography}

\end{document}
